% report.tex -- расчётно-пояснительная записка НИРС (магистратура),
% продолжение ВКР бакалавра И. П. Шаманова 2025 г.
% Сборка: latexmk -pdf -interaction=nonstopmode report.tex
\documentclass[12pt,a4paper]{article}
\usepackage[T2A]{fontenc}
\usepackage[utf8]{inputenc}
\usepackage{cmap}
\usepackage[russian]{babel}
\usepackage{amsmath,amssymb}
\usepackage{booktabs}
\usepackage{multirow}
\usepackage{array}
\usepackage{geometry}
\usepackage{enumitem}
\usepackage{pgfplots}
\usepackage{microtype}
\usepackage[htt]{hyphenat}
\usepackage{hyperref}
\setlength{\emergencystretch}{3.5em}
\tolerance=2000
\pgfplotsset{compat=1.17,
  every axis/.append style={label style={font=\small},
    tick label style={font=\footnotesize}}}
\geometry{left=2.5cm,right=2cm,top=2cm,bottom=2.5cm}
\hypersetup{colorlinks=true,linkcolor=black,urlcolor=blue,citecolor=black}

\newcommand{\divg}{\operatorname{div}}
\newcommand{\B}{\mathbf{B}}
\newcommand{\vv}{\mathbf{v}}
\newcommand{\code}[1]{\texttt{#1}}
\newcommand{\Rl}{\textsuperscript{\scriptsize Р}}   % реализовано
\newcommand{\Pv}{\textsuperscript{\scriptsize П}}   % проверено автотестом
\newcommand{\Iz}{\textsuperscript{\scriptsize И}}    % измерено (единичная диагностика)
\newcommand{\Gp}{\textsuperscript{\scriptsize Г}}   % гипотеза
\newcommand{\Pl}{\textsuperscript{\scriptsize Пл}}  % план

\title{Развитие и верификация программного комплекса для решения\\
задач идеальной магнитной гидродинамики: блочно-структурированный\\
решатель второго порядка и общий контур проверки достоверности\\[6pt]
\large Расчётно-пояснительная записка к научно-исследовательской работе}
\author{И.\,П.~Шаманов\\[2pt]
\normalsize научный руководитель В.\,В.~Лукин\\
\normalsize МГТУ им. Н.\,Э.~Баумана, кафедра <<Прикладная математика>>}
\date{4 сентября 2026 г.}

\begin{document}
\maketitle

\begin{abstract}
\noindent
Работа продолжает выпускную квалификационную работу бакалавра
\cite{shamanov2025}, где был построен решатель идеальной магнитной
гидродинамики первого порядка на неструктурированной треугольной сетке по
схеме Авдеевой--Лукина \cite{avdeeva2019,galanin2007}. Сама та работа называет
главным своим ограничением высокую диссипацию метода и рекомендует перейти к
более высокому порядку. Здесь это сделано: построен решатель второго порядка на
декартовой блочно-структурированной сетке с адаптивным измельчением, MUSCL,
SSP-RK2 и constrained transport, с распараллеливанием MPI и OpenMP.

Прежде чем сравнивать, старую схему пришлось починить. Контролируемое
воспроизведение на зафиксированной ревизии исходного кода показало, что она
расходится к отрицательным плотности и давлению на всех разрывных тестах.
Создан исправленный профиль с реестром численных дельт и регрессионными
тестами; все канонические задачи доведены до конца, физически допустимы,
консервативны и дискретно бездивергентны.

Сравнение ведётся по общей метрике при равном разрешении и против эталона,
построенного независимой схемой. При $N=400$ новая схема разрешает контактный
разрыв в задаче Брио--Ву в 3.9 ячейки против 36.8 у старой --- в 9.4 раза
резче, --- а ошибка при $N\geq256$ меньше в 17--19 раз. Выигрыш раскладывается
на две части: 5.0 раза даёт переход на декартову сетку при неизменном первом
порядке, ещё 3.4--3.6 --- второй порядок. На тесте с переносом петли поля, где
выпускная работа и констатирует высокую диссипацию, за два прохода области
старая схема сохраняет $31\,\%$ магнитной энергии, новая --- $84\,\%$.
Дискретная дивергенция поля у обеих остаётся на уровне ошибок округления.

Не заявляется до прохождения соответствующих проверок: готовность
GPU-версии (сборка невозможна на доступном оборудовании, выполнен аудит
переносимости) и многоузловое масштабирование (нет доступа к кластеру; на одном
узле измерены сильное и слабое). Полная трассировка утверждений к сырым данным
приведена в \code{docs/CLAIM\_TO\_EVIDENCE.md} той же ревизии репозитория.
\end{abstract}

\medskip
\noindent\textit{Обозначения статуса результата:}
\Rl~--- реализовано в исходном коде;
\Pv~--- проверено автоматическим тестом;
\Iz~--- измерено единичной диагностикой (не benchmark);
\Gp~--- исследовательская гипотеза;
\Pl~--- запланировано.

\tableofcontents
\newpage

% =====================================================================
\section{Введение}

\subsection{Зачем считать магнитную гидродинамику}

Магнитная гидродинамика описывает проводящую среду — плазму, ионизованный газ,
жидкий металл — вместе с магнитным полем, которое эта среда за собой увлекает и
которое на неё же давит. Обратная связь работает в обе стороны: поле нельзя
задать заранее, оно часть решения. Из-за этого почти всякая содержательная
постановка решается численно.

Астрофизика даёт самые наглядные примеры. В тесной двойной системе вещество
перетекает с одной звезды на другую и закручивается в аккреционный диск; если у
принимающей звезды сильное собственное поле, диск разбивается на зоны, между
которыми возникают токовые слои, и картина течения меняется качественно
\cite{zhilkin2010,zhilkin2012}. Солнечная вспышка — это перезамыкание силовых
линий, при котором энергия поля за минуты переходит в кинетическую и тепловую.
В установках управляемого термоядерного синтеза вопрос удержания плазмы —
это в первую очередь вопрос об устойчивости магнитной конфигурации. Во всех
трёх случаях интересное происходит на масштабах много меньших, чем размер
области: токовый слой тоньше диска на порядки, фронт ударной волны занимает
доли процента домена. Это и есть довод в пользу адаптивного измельчения сетки —
считать везде с шагом, нужным в самом тонком месте, слишком дорого.

\subsection{Что ломается в расчёте}
\label{sec:intro-div}

У магнитного поля нет источников: силовые линии нигде не начинаются и не
кончаются. Записанное как $\divg\B=0$, это условие в непрерывных уравнениях
выполняется само собой — если поле было соленоидальным вначале, уравнение
индукции его таким и оставит.

У разностной схемы такой гарантии нет. Ошибки аппроксимации накапливаются, и в
расчёте заводится то, чего в природе не бывает, — «магнитный заряд». Он не
безобиден: в разностных уравнениях появляется сила, действующая вдоль поля и
разгоняющая вещество там, где разгонять нечему. Дальше — искажение решения,
отрицательное давление и остановка счёта. Задача, таким образом, не в том, чтобы
сделать невязку $\divg\B$ малой, а в том, чтобы схема не производила её вовсе.

Ответов на это два. Первый — разрешить дивергенции появляться и затем убирать
её: сносить источниковыми членами (восьмиволновой метод) или вычитать
градиентную часть, решая уравнение Пуассона. Так устроены, в частности, расчёты
течений в тесных двойных системах у Жилкина и Бисикало \cite{zhilkin2010}. Второй
ответ — хранить магнитное поле не в центрах ячеек, а на их гранях, и обновлять
его циркуляцией электрического поля по контуру. Тогда дискретная дивергенция
сокращается тождественно и остаётся нулём с точностью до ошибок округления, без
всякой чистки. Схема Авдеевой--Лукина \cite{avdeeva2019,galanin2007},
унаследованная из выпускной работы бакалавра, и constrained transport нового
решателя принадлежат этой второй ветви. Плата за неё — усложнённое хранение
поля и необходимость согласовывать граневое и клеточное представления.

\subsection{Что даёт эта работа}

Коротко: раньше был решатель первого порядка на неструктурированной треугольной
сетке — бездивергентный, но слишком диссипативный и без адаптации. Здесь
построен решатель на декартовой блочно-структурированной сетке с адаптивным
измельчением, второго порядка по пространству и времени, с распараллеливанием
MPI и OpenMP. Результат: фронты резче, $\divg\B$ по-прежнему нулевая, появилась
адаптация.

Насколько резче — измерено. Тест с переносом петли магнитного поля построен так,
что идеальная схема сохранила бы магнитную энергию петли целиком; выпускная
работа именно на нём констатирует «свойство высокой диссипации выбранного
метода» и рекомендует перейти к схеме более высокого порядка. При одинаковой
сетке $128\times64$ за два прохода области старая схема сохраняет $31\,\%$
магнитной энергии, новая — $84\,\%$. На задаче Брио--Ву при одинаковой сетке
$N=400$ контактный разрыв, который старая схема размазывает на 37 ячеек, новая
разрешает в 4 (\S\ref{sec:res-rp3}).

Работа делает и третье, чего не планировалось. Контролируемое воспроизведение
исторического решателя показало, что на зафиксированной ревизии он расходится:
на всех разрывных тестах счёт уходит в нефизичное состояние. Поэтому старую
схему пришлось сначала починить — иначе сравнивать было бы не с чем.

\subsection{Состояние исторического решателя}

В выпускной работе бакалавра \cite{shamanov2025} реализован комплекс решения
двумерных и осесимметричных задач идеальной МГД на неструктурированных
треугольных сетках. Схема расщепляется по физическим процессам: газодинамические
потоки на рёбрах считает приближённый решатель задачи Римана (HLL, HLLC-L,
HLLD), нормальные компоненты поля на рёбрах обновляются законом Фарадея через
теорему Стокса, а поле в центрах ячеек восстанавливается базисными функциями
Равьяра--Тома. Порядок первый и по времени, и по пространству.

Воспроизведение выполнялось на официально зафиксированной ревизии исходного
кода (\code{9d0f60e}, репозиторий \code{Ship-Vano/MHD2D}); каждый вход и выход
помещался в изолированный каталог с фиксацией контрольных сумм. Установлено:

\begin{enumerate}[nosep]
\item на задаче Брио--Ву решатель расходится к нефизичному состоянию
  ($\rho_{\min}=-0.177$, $p_{\min}=-0.073$) на итерации~68;
\item на задачах о циркулярно-поляризованной альфвеновской волне и о петле
  магнитного поля нефизичное состояние возникает уже на первом обновлении;
\item во всех трёх случаях процесс завершается с нулевым кодом возврата, потому
  что контроль ограничен проверкой на \code{NaN}.
\end{enumerate}

Аудит исходного кода объясняет, откуда это берётся. Максимум $\divg\B$ оценивался
со знаком, а не по модулю, и потому скрывал отрицательную по модулю ошибку.
Отката потока и счётчиков нефизичных состояний не было вовсе. Число Куранта
задано в коде жёстко и не совпадает с указанным в тексте работы. Формулы правой
скорости волны в HLLD расходятся между CPU- и CUDA-веткой. У угловых
треугольников терялся один из двух граничных углов при построении отображения
<<граничный элемент~$\to$~фиктивный элемент>>.

\subsection{Цель и задачи}

Цель — получить два независимо собираемых и проверяемых решателя идеальной МГД с
общим контуром проверки достоверности и выяснить, что даёт переход ко второму
порядку и адаптивной сетке.

Задачи:
\begin{enumerate}[nosep]
\item восстановить исторический решатель и устранить выявленные дефекты в
  отдельном профиле с полным реестром численных дельт и регрессионными тестами;
  довести канонические задачи до конца и подтвердить физическую допустимость,
  консервативность и бездивергентность;
\item разработать блочно-структурированный решатель второго порядка на
  декартовых сетках с адаптивным измельчением: MUSCL, SSP-RK2, constrained
  transport, HLLD с положительность-сохраняющим откатом;
\item построить воспроизводимый контур проверки достоверности: автоматические
  тесты ядра и канонических задач, манифесты запусков, машинно вычисляемые
  метрики;
\item сравнить старую и новую схемы на общих постановках по общей метрике и
  сопоставить результаты с данными выпускной работы \cite{shamanov2025}, статьи
  Авдеевой--Лукина \cite{avdeeva2019} и каталога тестов Athena;
\item охарактеризовать стоимость расчёта и параллельную эффективность.
\end{enumerate}

Проверяемая гипотеза~H1: на рассматриваемом классе задач схема второго порядка
на блочно-структурированной сетке с адаптивным измельчением способна дать лучшее
соотношение «ошибка — время — память», чем схема первого порядка на
неструктурированной сетке. Гипотеза не принимается заранее; её частичное или
полное опровержение допустимо и было бы результатом.

% =====================================================================
\section{Математическая постановка}

Идеальная МГД — это уравнения газовой динамики, дополненные магнитным полем и
записанные в предположении бесконечной проводимости среды: поле вморожено в
вещество и переносится вместе с ним. Вязкость, сопротивление и теплопроводность
не учитываются. Все четыре уравнения ниже — законы сохранения: массы, импульса,
энергии и магнитного потока. Дивергентная форма выбрана не для красоты записи, а
потому что конечно-объёмная схема аппроксимирует именно её и наследует
сохранение автоматически.

В определение вектора магнитной индукции включён множитель $1/\sqrt{4\pi}$,
чтобы в уравнениях не появлялись явные $4\pi$:
\begin{align}
 \partial_t\rho + \nabla\cdot(\rho\vv) &= 0, \label{eq:mass}\\
 \partial_t(\rho\vv)+\nabla\cdot\!\left(\rho\vv\vv + p_T\hat I - \B\B\right) &= 0,
   \label{eq:mom}\\
 \partial_t E+\nabla\cdot\!\left((E+p_T)\vv - \B(\vv\cdot\B)\right) &= 0,
   \label{eq:ene}\\
 \partial_t\B + \nabla\cdot(\vv\B - \B\vv) &= 0, \qquad \divg\B = 0.
   \label{eq:faraday}
\end{align}
Здесь $\rho$ --- плотность, $\vv=(u,v,w)^{\mathsf T}$ --- скорость,
$\B=(B_x,B_y,B_z)^{\mathsf T}$ --- магнитная индукция,
$E=\rho\varepsilon+\tfrac12\rho|\vv|^2+\tfrac12|\B|^2$ --- полная энергия
единицы объёма, $p_T=p+\tfrac12|\B|^2$ --- полное давление,
$p=(\gamma-1)\bigl(E-\tfrac12\rho|\vv|^2-\tfrac12|\B|^2\bigr)$,
$\gamma$ --- показатель адиабаты. Все величины зависят от $(x,y,t)$, но скорость
и поле сохраняют три компоненты: это режим 2.5D, в котором поле может
поворачиваться из плоскости расчёта. Без третьей компоненты не воспроизвести
вращательные разрывы, а с ними и задачу Брио--Ву.

Магнитное поле входит в уравнение импульса двумя слагаемыми, и оба имеют
физический смысл: $\tfrac12|\B|^2$ — магнитное давление, добавляющееся к
газовому, а $-\B\B$ — натяжение силовых линий, из-за которого поле сопротивляется
изгибу. Второе слагаемое и делает МГД непохожей на газовую динамику: среда
приобретает выделенное направление.

Если применить оператор дивергенции к~\eqref{eq:faraday}, получится
$\partial_t(\divg\B)=0$. Соленоидальность не нужно поддерживать — она сохраняется
сама, если выполнена в начальный момент. Именно это свойство и теряет разностная
схема (\S\ref{sec:intro-div}), и именно его восстанавливает constrained transport
(\S\ref{sec:ct}).

Скорости, с которыми возмущения бегут по такой среде, — собственные значения
матрицы Якоби потока. В направлении~$x$ их семь:
\[
 \lambda_{4}=u,\quad
 \lambda_{2,6}=u\mp c_a,\quad
 \lambda_{1,7}=u\mp c_f,\quad
 \lambda_{3,5}=u\mp c_s,
\]
\[
 c_a=\frac{|B_x|}{\sqrt\rho},\qquad
 c_{f,s}=\left\{
 \frac{\gamma p+|\B|^2\pm\sqrt{(\gamma p+|\B|^2)^2-4\gamma p\,B_x^2}}{2\rho}
 \right\}^{1/2}.
\]
Одна волна переносится вместе с веществом, две альфвеновские бегут вдоль поля со
скоростью $c_a$, а четыре магнитозвуковые — быстрые и медленные — со скоростями
$c_f$ и $c_s$. Эту семёрку и приближает решатель HLLD (\S\ref{sec:hlld}).

Существенно, что собственные значения могут совпадать: при $B_x\to0$ медленная
волна сливается с альфвеновской, а при $\B\to0$ обе — с контактной. Система не
строго гиперболична, а поток не выпукл \cite{brio1988}. Практическое следствие
неприятное: в решении задачи о распаде разрыва могут появляться волны, которых в
газовой динамике не бывает, — составные, где ударная волна срослась с
разрежением. Классические решатели, разлагающие решение по семи регулярным
волнам, такую конфигурацию не представляют, и потому для теста Брио--Ву не
существует пригодного точного решения (\S\ref{sec:res-briowu-t06}).

% =====================================================================
\section{Исторический решатель и его исправление}
\label{sec:legacy}

\subsection{Схема первого порядка (профиль \code{legacy\_vkr})}

Профиль \code{legacy\_vkr} сохраняет метод ВКР без изменения математики:
конечный объём на неструктурированной треугольной сетке, первый порядок по
пространству и времени, HLLD-поток, нормальная компонента $B_n$ на рёбрах,
узловая ЭДС как арифметическое среднее тангенциальных компонент HLLD-потоков
примыкающих рёбер, реконструкция Равьяра--Тома. Число Куранта и параметры
задаются версионируемым описанием задачи. Опубликованные результаты не
изменяются задним числом.

Контролируемое воспроизведение (см.~введение) зафиксировано отдельным
протоколом: расхождение на всех разрывных тестах есть свойство именно этой
ревизии, а не ошибка запуска.

\subsection{Профиль \code{legacy\_corrected} и реестр численных дельт}

Исправления реализованы как версионируемая заплатка к неизменяемому исходному
коду и собраны в профиль \code{legacy\_corrected}. Каждая дельта имеет
регрессионный тест (набор \code{legacy\_corrected\_kernel}, проходит).

\begin{center}\footnotesize
\begin{tabular}{p{4.6cm}p{10.4cm}}
\toprule
Дельта & Причина и физический смысл \\
\midrule
Исправленный HLLD-поток и наблюдаемый откат на HLL & промежуточные состояния и
сигнал-скорости проверяются на конечность и положительность; недопустимое
промежуточное состояние не должно распространять \code{NaN}, а заменяется
положительность-сохраняющим потоком HLL со счётчиком срабатываний\Rl\Pv \\
Восстановление примитивных/консервативных переменных & давление должно вычитать
кинетическую и магнитную энергию ровно один раз; проверено круговым
преобразованием $(\rho,u,v,w,p)$\Pv \\
Число Куранта для произвольного треугольника & шаг ограничивается величиной
$\mathrm{CFL}\cdot A_K/\sum_e \lambda_e \ell_e$ (площадь ячейки, спектральный
радиус и длина ребра); скрытый нижний порог шага удалён\Rl\Pv \\
Отображение <<граница~$\to$~фиктивный элемент>> & у углового треугольника два
граничных ребра; хранится взаимно-однозначное соответствие
ребро$\leftrightarrow$фиктивный элемент\Rl\Pv \\
Оценка $\divg\B$ и ориентация CT & знаковый максимум скрывал отрицательную по
модулю ошибку; используется $\max|\cdot|$ и нормированная норма; знак закона
Фарадея на ребре выражен явно через касательный вектор\Rl\Pv \\
Реконструкция Равьяра--Тома & поле в центре восстанавливается из
знако-ориентированной относительно элемента нормальной компоненты $B_n$, а не
из порядка узлов; постоянное поле восстанавливается точно на неправильных
треугольниках\Pv \\
Энергия после CT & подгонка $E$ после реконструкции поля искусственно меняла
полную энергию; исправленный путь сохраняет консервативную $E$ от потока, а
изменение магнитной энергии от реконструкции логируется отдельно\Rl \\
Диагностика физического отказа & выход с нулевым кодом не отличает физический
отказ; при недопустимом состоянии выводится дамп (шаг, время, элемент,
$\rho$, $p$) и процесс завершается ненулевым кодом\Rl\Pv \\
\bottomrule
\end{tabular}
\end{center}

Оставшийся вне рамок настоящей работы пробел: типы задач <<вращающийся
цилиндр>> и <<вихрь Орзага--Танга>> ещё используют жёстко заданное в коде число
Куранта; исправление <<CFL --- явный параметр конфигурации>> применено к
задачам Брио--Ву, CP-альфвен и петле поля.

\subsection{Сохранённый метод второго порядка не активирован}

В исходном коде присутствует функция ограничителя \code{minmod}, однако
MUSCL-реконструкция \emph{не вызывается}: профиль \code{legacy\_corrected}
намеренно остаётся первопорядковым, чтобы служить честной базой сравнения для
схемы второго порядка (п.~\ref{sec:new}).

% =====================================================================
\section{Новый блочно-структурированный решатель}
\label{sec:new}

\subsection{Что и где хранится}

Решатель работает на декартовой блочно-структурированной сетке, которая умеет
измельчаться там, где это нужно решению; вся работа с иерархией уровней делегирована
библиотеке AMReX~\cite{amrex}. Схема наследует у метода Авдеевой--Лукина главную
идею — разное хранение для газовых величин и для магнитного поля, — но переносит
её на прямоугольную ячейку.

Плотность, три компоненты импульса, полная энергия и $B_z$ живут в центрах ячеек:
это величины, которые переносятся потоками через грани. Плоскостное поле хранится
иначе — $B_x$ на $x$-гранях, $B_y$ на $y$-гранях, — и именно эти граневые значения
считаются первичными данными. Электрическое поле $E_z$ определено в узлах сетки.
Клеточные $B_x$ и $B_y$, нужные решателю Римана и выводу результатов, получаются
полусуммой противоположных граней — на прямоугольнике это и есть реконструкция
Равьяра--Тома.

Такое разнесение выглядит избыточным, но именно оно даёт бездивергентность
даром (\S\ref{sec:ct}). Плата — необходимость держать поле в двух
представлениях согласованными.

\subsection{Реконструкция и решатель Римана}
\label{sec:hlld}

Схема первого порядка считает величину постоянной внутри ячейки, поэтому размывает
любой разрыв на десяток ячеек. MUSCL-реконструкция вместо этого проводит внутри
ячейки прямую и берёт на грань её значение на краю. Наклон прямой ограничивается,
иначе на разрыве появятся осцилляции; ограничитель (minmod, van~Leer или MC)
задаётся конфигурацией, а его выбор измерен в~\S\ref{sec:res-rp1}. Нормальная
компонента поля на грань не реконструируется: она берётся из граневого массива,
иначе рассогласуется с бездивергентным обновлением.

Дальше на каждой грани решается задача о распаде разрыва. Точное её решение в МГД
дорого и не всегда существует в привычном виде, поэтому используется приближённый
решатель HLLD~\cite{miyoshi2005}. Физическая идея простая: веер волн,
разбегающихся от разрыва, приближается пятью волнами — двумя быстрыми
магнитозвуковыми по краям, двумя вращательными и контактной в середине, — а
область между ними разбивается на четыре промежуточных состояния с постоянными
значениями. Именно наличие вращательных волн отличает HLLD от более грубых HLL и
HLLC и позволяет не размазывать те разрывы, на которых поле поворачивается.

Крайние скорости веера оцениваются как
\[
 S_L=\min(u_L,u_R)-\max(c_{fL},c_{fR}),\qquad
 S_R=\max(u_L,u_R)+\max(c_{fL},c_{fR}),
\]
то есть так же, как в~\cite{avdeeva2019}. Оценка шире эйнфельдтовой, и потому
осреднённый по вееру поток HLL, построенный на ней, сохраняет положительность при
выполнении условия Куранта. Это существенно для отката, описанного ниже.

Промежуточные состояния HLLD вырождаются: знаменатель обращается в ноль при
$(S_\alpha-u_\alpha)(S_\alpha-S_M)\to B_x^2/\rho$ и вблизи вакуума. Вместо того
чтобы надеяться, что такие состояния не встретятся, решатель их проверяет.
Результат HLLD принимается, только если он конечен, а плотности осреднённого и
звёздных состояний положительны; иначе поток заменяется потоком HLL, а
срабатывание попадает в счётчик. Счётчик здесь не декоративен: тесты падают, если
откат сработал на гладкой задаче, где его быть не должно. На допустимом состоянии
поток HLLD от проверки не изменился побитово\Pv.

\subsection{Почему поле обновляется отдельно}
\label{sec:ct}

Если обновлять $B_x$ и $B_y$ так же, как плотность, — балансом потоков через
грани ячейки, — дискретная дивергенция начнёт расти, и дальше происходит то, что
описано во введении. Constrained transport обходит это, обновляя не клеточные
значения, а граневые, и не потоками, а циркуляцией электрического поля по контуру
грани.

Устроено это так. Поток HLLD через $x$-грань уже содержит ЭДС: его компонента при
$B_y$ равна $-E_z$; на $y$-грани компонента при $B_x$ равна $+E_z$. В каждом узле
сетки сходятся четыре такие грани, и их ЭДС усредняются в одно узловое значение
$E_z$. Затем $B_x$ и $B_y$ на гранях меняются по дискретному закону Фарадея —
разности узловых $E_z$ вдоль грани.

Выигрыш виден, если выписать дискретную дивергенцию ячейки: в неё входят четыре
граневых приращения, и узловые $E_z$ в этой сумме сокращаются попарно. Дивергенция
не уменьшается численно — она не меняется вообще, с точностью до ошибок
округления\Pv. Соленоидальность становится тождеством схемы, а не результатом,
за которым надо следить.

Способ усреднения ЭДС в узле выбирается конфигурацией: простое среднее четырёх
граней (Balsara--Spicer~\cite{balsara1999}) либо форма
$E_z=2\langle E_z\rangle_{\text{грани}}-\langle E_z\rangle_{\text{ячейки}}$
(Gardiner--Stone~\cite{gardiner2005}), добавляющая диссипацию. Второй вариант
принят по умолчанию: он уменьшает перелёты у фронтов в полтора раза без потери
точности (\S\ref{sec:res-rp1}).

\subsection{Интегратор по времени}

По времени схема двухстадийная — SSP-RK2, он же метод Хойна:
\[
 U^{(1)}=U^n+\Delta t\,L(U^n),\qquad
 U^{n+1}=\tfrac12 U^n+\tfrac12\!\left(U^{(1)}+\Delta t\,L(U^{(1)})\right).
\]
Обе стадии — шаги Эйлера, а итог их выпуклая комбинация с положительными
весами. Отсюда и название: если шаг Эйлера при допустимом числе Куранта не
увеличивает полную вариацию, то её не увеличивает и весь шаг. В литературе эта
формула известна как TVD RK2 Шу--Ошера~\cite{shu1988,gottlieb2001}; терминология
здесь согласована намеренно, поскольку в тексте ВКР та же схема названа иначе.
Для линейного оператора $L(y)=\lambda y$ множитель перехода равен
$1+a+\tfrac12 a^2$ при $a=\lambda\Delta t$ — ровно два первых члена ряда для
$\mathrm e^{a}$, что и проверяет автотест\Pv. Каждая стадия делает своё
CT-обновление, а выпуклая комбинация бездивергентных полей бездивергентна, так
что интегратор соленоидальность не портит.

Про границы этого обоснования стоит сказать прямо, потому что измерение его не
подтвердило. TVD-свойство — теоретическая гарантия, и на задаче Брио--Ву она не
даёт ничего измеримого: SSP-RK2 и не-SSP схема средней точки совпадают по
относительной $L_1$ до четвёртой значащей цифры, а по амплитуде перелёта
расходятся на $0.1\,\%$; проверено на всём допустимом шаге вплоть до
$\mathrm{CFL}=0.9$ (\S\ref{sec:res-rp1}). TVD-RK3 при в полтора раза большей
стоимости тоже не измеряется лучше. Что действительно меняет картину — это сам
второй порядок по времени: явный Эйлер при почти той же $L_1$ удваивает перелёт.
Поэтому SSP-RK2 выбран за гарантию и цену, а не за наблюдаемый выигрыш.

\subsection{Адаптивное измельчение сетки}

Иерархия уровней строится стандартными средствами AMReX: фантомные ячейки и
фантомные грани заполняются интерполяцией с грубого уровня, причём для
граневого поля используется бездивергентная интерполяция — обычная разрушила бы
то, что даёт constrained transport. Ячейки помечаются к измельчению по
относительному градиенту плотности, по плотности тока $j_z=(\nabla\times\B)_z$
или по обоим критериям сразу. Подциклирования по времени нет: все уровни идут
общим шагом $\Delta t$.

Есть место, где такая иерархия ломает консервативность, и его пришлось чинить
отдельно. Грубая ячейка у стыка уровней обновляется грубым потоком, а мелкие
ячейки за той же гранью — набором мелких потоков; эти два числа не совпадают.
Разница накапливается, и суммарная масса перестаёт сохраняться. Лечится
регистром потоков, который запоминает оба вклада и поправляет грубую ячейку на
их разность; измерение — в~\S\ref{sec:res-amr}, где дрейф массы падает с
$7\!\cdot\!10^{-5}$ до уровня округления.

Расчёт ведётся на CPU. Заголовки вычислительных ядер помечены
device-квалификаторами и не зависят от AMReX-контейнеров, то есть к переносу на
GPU готовы, но горячие циклы остаются хостовыми; что именно мешает и что нужно
сделать — в~\S\ref{sec:res-gpu}.

% =====================================================================
\section{Методика проверки достоверности}
\label{sec:vv}

\subsection{Что значит «результат воспроизводим»}

Расчёт, который нельзя повторить, результатом не является, поэтому каждый прогон
сопровождается манифестом: ревизия репозитория и признак незакоммиченных
изменений, компилятор и флаги, версии зависимостей, контрольная сумма
конфигурации, платформа. Конфигурация задачи читается из JSON и проверяется
строго — неизвестный ключ приводит к отказу до того, как выделена память под
сетку, так что опечатка в имени параметра не проходит молча. Предусмотрен режим
без служебного ввода-вывода: замеры не должны включать запись файлов
результатов.

Исторический решатель воспроизводится сложнее, потому что его исходный код
неизменяем по договорённости. Сценарий запуска клонирует зафиксированную
ревизию, накладывает заплатку с предварительной проверкой, собирает, прогоняет
тесты и формирует манифест — в том числе когда контроль качества не пройден.
Неудачный прогон сохраняется наравне с удачным: отказ тоже результат, и он
должен быть воспроизводим. Сетка строится своя, структурная треугольная (два
прямоугольных треугольника на прямоугольник), чтобы не зависеть от внешнего
генератора; её хэш входит в манифест.

\subsection{Автоматические тесты}

Набор CTest содержит 21 тест и проходит целиком. Тесты делятся на три группы по
тому, что именно они защищают.

\emph{Ядро схемы.} Круговое преобразование переменных на 2000 случайных
состояниях; предельные случаи быстрой магнитозвуковой скорости; согласованность
HLLD, то есть $\mathbf F(q,q)=\mathbf F_{\text{физ}}(q)$, на 3000 состояниях;
сверхзвуковые ветви; вырождения ($B_n=0$, состояние Брио--Ву, вращательный
разрыв); 20000 состояний в широком диапазоне параметров на конечность потока;
детерминированные срабатывания отката HLLD$\to$HLL; инварианты ограничителя
(сохранение постоянной, обнуление на экстремуме, нечётность, TVD-оценка);
узловая ЭДС; коэффициенты и порядок SSP-RK2.

\emph{Канонические задачи целиком.} Постоянное состояние (схема не должна его
портить), вихрь Орзага--Танга, вращающийся цилиндр и МГД-взрыв — с проверкой
итоговых диапазонов $\rho$ и $p$ против литературных. Такой тест ловит не только
падение расчёта, но и смещение фронта или потерю признака решения. Наблюдаемый
порядок на циркулярно-поляризованной альфвеновской волне проверяется с порогом
$\geq1.8$.

\emph{Свойства, которые легко потерять незаметно.} Консервативность на стыке
уровней и при перестроении сетки; независимость результата от разбиения по
MPI-рангам; воспроизводимость манифеста; отсутствие структурной зависимости
слоя ядер от AMReX; согласие с независимым эталоном Брио--Ву. Отдельно стоит
сказать про тест консервативности: он устроен так, чтобы падать при вырождении
постановки — если мелкий уровень покрывает весь периодический домен, стыка
уровней нет и проверять нечего, — и повторяет прогон с выключенным согласованием
потоков, требуя, чтобы дефект действительно проявился. Тест, который проходит
потому, что ничего не проверяет, хуже отсутствующего.

\subsection{С чем сравниваются результаты}

Опубликованные в \cite{shamanov2025,avdeeva2019} графики качественные:
оцифрованных данных нет. Поэтому сравнение с ними ведётся по признакам решения —
диапазоны величин, число и положение фронтов, рост $\divg\B$, наличие паразитных
осцилляций, — а не как $L$-ошибка относительно картинки.

Для задачи Брио--Ву этого мало, потому что она и есть главный тест сравнения
схем. Точное решение здесь недостижимо: конфигурация содержит составную волну,
которую классические семиволновые решатели не представляют. Поэтому эталон
построен независимой схемой — центральной схемой Куртганова--Тадмора, не
решающей задачу Римана вообще и не разделяющей с проверяемой схемой ни одной
строки кода (\S\ref{sec:res-briowu-t06}). Более ранняя версия работы
использовала в этой роли сошедшийся расчёт самого нового решателя; такой эталон
доказывает лишь самосогласованность, и он заменён.

% =====================================================================
\section{Результаты}
\label{sec:res}

\subsection{Исправленный решатель: сквозные расчёты}
\label{sec:res-legacy}

Все канонические задачи профиля \code{legacy\_corrected} доведены до конца на
структурной треугольной сетке (заплатка \code{7d7314c6} к ревизии
\code{9d0f60e}, компилятор GCC~15.2). Во всех случаях: состояние конечно и
положительно (страж положительности не срабатывал), нормированная норма
$\divg\B$ на уровне ошибок округления, законы сохранения массы, импульса и
энергии замыкаются с точностью до округления против интеграла граничного
потока, откат HLLD$\to$HLL ни разу не задействован.

\begin{center}\footnotesize\setlength{\tabcolsep}{4.5pt}
\begin{tabular}{lccccccc}
\toprule
Задача & Сетка & $t_{\text{кон}}$ & CFL & $N_{\text{ит}}$ & $\rho_{\min}$ &
норм. $\|\divg\B\|$ & Невязка сохр. \\
\midrule
Брио--Ву              & $128\times4$   & 0.1  & 0.1 & 4307  & 0.112 & $6.2\!\cdot\!10^{-15}$ & $4.0\!\cdot\!10^{-16}$ \\
CP-альфвен, $30^\circ$ & $32\times56$   & 1.0  & 0.1 & 2211  & 0.893 & $2.5\!\cdot\!10^{-15}$ & $1.8\!\cdot\!10^{-13}$ \\
Петля поля (масшт.)   & $64\times32$   & 2.0  & 0.1 & 10762 & 1.000 & $2.2\!\cdot\!10^{-15}$ & $1.9\!\cdot\!10^{-14}$ \\
Петля поля (Athena)   & $64\times37$   & 4.62 & 0.1 & 18159 & 1.000 & $3.2\!\cdot\!10^{-15}$ & $1.3\!\cdot\!10^{-13}$ \\
Вращающийся цилиндр   & $128\times128$ & 0.15 & 0.5 & 820   & 0.743 & $1.4\!\cdot\!10^{-14}$ & $2.1\!\cdot\!10^{-12}$ \\
\bottomrule
\end{tabular}
\end{center}
\noindent\footnotesize
Источник: \code{legacy\_corrected\_endtoend.json},
\code{LEGACY\_CORRECTED\_T05\_ENDTOEND.md} (каталоги \code{benchmarks/summary/}, \code{docs/});
строка «вращающийся цилиндр» --- прогон
\code{benchmarks/raw/legacy\_corrected/rotor\_128/} с манифестом и пройденным
quality gate. Для CP-альфвен усреднённая
относительная $L_1$-ошибка по $(v_\perp,v_z,B_\perp,B_z)$ равна $0.370$
(первый порядок); для петли поля отношение магнитной энергии к начальной
$\approx0.17$ (сильная диссипация).\normalsize

\paragraph{Вихрь Орзага--Танга: причина отказа и её устранение.}
Исходно расчёт на сетке $N=128$ \emph{не доводился до конца}: страж
положительности срабатывал на итерации~1729 при $t=0.43120$. Причина установлена по
сохранённому состоянию отказавшей ячейки (центр $(0.354,\,0.044)$):
\[
  \rho=0.3665,\quad
  \tfrac12\rho|\vv|^2=1.68\!\cdot\!10^{-3},\quad
  \tfrac12|\B|^2=0.26035,\quad
  e=0.26196 .
\]
Тепловая энергия получается как разность почти равных величин,
$e-\tfrac12\rho|\vv|^2-\tfrac12|\B|^2=-7.0\!\cdot\!10^{-5}$, то есть
составляет лишь $0.03\,\%$ полной энергии: ячейка магнитно доминирована
($\beta\to0$), и ошибка аппроксимации первого порядка в $e$ превышает само
тепловое слагаемое. Это \emph{катастрофическое сокращение} в
$p=(\gamma-1)\bigl(e-\tfrac12\rho|\vv|^2-\tfrac12|\B|^2\bigr)$, а не дефект
начальных данных: начальная невязка магнитного потока равна
$1.0\!\cdot\!10^{-17}$, и переход к заданию $B_n$ через разности векторного
потенциала (\S\ref{sec:legacy}) точку отказа не сдвинул. Устранение
требует более подробной сетки ($N\approx400\text{--}800$, как в
\cite{shamanov2025,avdeeva2019}), схемы второго порядка либо формализма двух
энергий.

Источник накопления найден в шаге RT0. На каждой итерации клеточное поле $\B$
пересобирается из граневого, которым владеет CT, и магнитная энергия ячейки при
этом меняется. Полная энергия $E$ намеренно не корректировалась, чтобы потоковый
баланс оставался замкнутым, --- значит, всю эту разницу поглощает внутренняя
энергия, и в ячейке с $\beta\to0$ накопление уводит её в минус. Но
реконструкция --- это смена представления поля, а не нагрев газа, и менять $p$
она не должна. Если сдвигать $E$ на изменение магнитной энергии
(\code{ctEnergyMode: preserve\_internal}, решение~D-009), полная энергия меняется
ровно на уже логируемую величину, а отказ исчезает:

\begin{center}\footnotesize
\begin{tabular}{lcc}
\toprule
Режим & Итераций & Срабатываний пола давления \\
\midrule
\code{conservative} (прежний) + пол & 2017 & 14300 \\
\code{preserve\_internal} (по умолчанию) & 2012 & \textbf{0} \\
\bottomrule
\end{tabular}
\end{center}
\noindent\footnotesize
$128\times128$, CFL 0.5, $t=0.5$. Прежний режим сохранён только как измерение\Iz.\normalsize

\medskip
Структура воспроизводится (центральный ромб, симметричные уплотнения,
диагональные фронты), а диапазоны сужены диффузией первого порядка. Что это
именно вопрос разрешения, а не корректности, показано измерением: при
измельчении сетки все четыре экстремума движутся к литературным.

\begin{center}\footnotesize
\begin{tabular}{lcccc}
\toprule
& $\rho_{\min}$ & $\rho_{\max}$ & $p_{\min}$ & $p_{\max}$ \\
\midrule
$N=128$                       & 0.0953 & 0.4018 & 0.0417 & 0.4454 \\
$N=256$                       & 0.0886 & 0.4262 & 0.0379 & 0.4619 \\
Авдеева--Лукин, $N=800$       & 0.0870 & 0.4890 & 0.0280 & 0.4900 \\
\midrule
отклонение при $N=128$        & $+9.5\,\%$ & $-17.8\,\%$ & $+48.9\,\%$ & $-9.1\,\%$ \\
отклонение при $N=256$        & $\mathbf{+1.8\,\%}$ & $-12.8\,\%$ & $+35.4\,\%$ & $-5.7\,\%$ \\
\bottomrule
\end{tabular}
\end{center}
\noindent\footnotesize
Манифесты: \code{benchmarks/raw/legacy\_corrected/orszag\_tang\_\{128,256\}/},
оба с пройденным quality gate\Iz.\normalsize

\medskip
Расхождение одностороннее и ровно в ту сторону, которую предсказывает диффузия:
пики занижены, впадины завышены. При $N=256$ (4055 итераций) $\rho_{\min}$
отличается от эталонного на $1.8\,\%$, и пол по давлению по-прежнему не
срабатывает ни разу.

Исторический профиль \code{legacy\_vkr} на тех же задачах даёт \code{NaN}
(п.~\ref{sec:legacy}); исправления \emph{восстановили} физическую картину, а не
исказили её.

\subsection{Новый решатель: канонические задачи}
\label{sec:res-new}

Автономный драйвер \code{mhd2d\_verify} использует те же заголовки ядер, что и
основной решатель. Во всех расчётах ниже: $\max|\divg\B|$ на уровне округления,
откат HLLD$\to$HLL не задействован, отрицательных ячеек нет.

\paragraph{Задача Брио--Ву.}\label{sec:res-briowu}
$\gamma=2$, $B_x=0.75$, $B_y=\pm1$, $t=0.1$. Профиль воспроизводит каноническую
\emph{составную (нерегулярную)} ветвь \cite{brio1988,avdeeva2019}: волна
разрежения быстрого семейства, составная волна со сменой знака $B_y$ вблизи
$x\approx0.5$, плато $\rho\approx0.5$, контактный разрыв, медленная ударная
волна (рис.~\ref{fig:briowu}).

\begin{figure}[htbp]\centering
\begin{tikzpicture}
\begin{axis}[width=0.49\textwidth,height=5cm,xlabel=$x$,ylabel=$\rho$,
  xmin=0,xmax=1,legend style={font=\tiny,at={(0.02,0.02)},anchor=south west},
  title={\footnotesize плотность}]
\addplot[gray,thick] table[x=x,y=rho] {figures/data/briowu1d_ref.dat};
\addplot[blue,mark=none,densely dashed] table[x=x,y=rho] {figures/data/briowu1d_legacy.dat};
\addplot[red,mark=none] table[x=x,y=rho] {figures/data/briowu1d_n3.dat};
\legend{эталон $N{=}2048$,legacy\_corrected,AMReX N3}
\end{axis}
\begin{axis}[xshift=0.5\textwidth,width=0.49\textwidth,height=5cm,xlabel=$x$,
  ylabel=$B_y$,xmin=0,xmax=1,title={\footnotesize $B_y$}]
\addplot[gray,thick] table[x=x,y=By] {figures/data/briowu1d_ref.dat};
\addplot[blue,mark=none,densely dashed] table[x=x,y=By] {figures/data/briowu1d_legacy.dat};
\addplot[red,mark=none] table[x=x,y=By] {figures/data/briowu1d_n3.dat};
\end{axis}
\end{tikzpicture}
\caption{Брио--Ву при равной сетке $N=400$: новая схема разрешает контактный
разрыв ($x\approx0.56$) в 3.9~ячейки против 36.8 у первопорядковой на
треугольниках --- в 9.4~раза резче. Смена знака $B_y$ вблизи $x\approx0.5$ ---
составная волна; её воспроизводят обе схемы, различаясь только тем, на сколько
ячеек она размазана. $t=0.1$, CFL~0.1; эталон --- независимая центральная схема
Куртганова--Тадмора при $N=6400$.
Источник: \code{docs/figures/data/briowu1d\_*.dat}.}
\label{fig:briowu}
\end{figure}

\paragraph{Циркулярно-поляризованная альфвеновская волна.}
$30^\circ$-геометрия Тота \cite{toth2000}, точное нелинейное решение, период
$t=1$. Наблюдаемый порядок сходимости по относительной $L_1$-ошибке $B_\perp$
близок ко второму (рис.~\ref{fig:conv}): $1.77\to1.79\to1.94$ на
последовательности $N=16,32,64,128$. Для сравнения, первопорядковая схема ВКР
на этой задаче даёт порядок $\approx1.2\text{--}1.3$ (табл.~3
в~\cite{shamanov2025}); при $N=128$ новая схема точнее примерно в 56~раз.

\begin{figure}[htbp]\centering
\begin{tikzpicture}
\begin{loglogaxis}[width=0.6\textwidth,height=5.2cm,xlabel={$N$},
  ylabel={относительная $L_1$-ошибка $B_\perp$},
  xtick={16,32,64,128},xticklabels={16,32,64,128},
  legend style={font=\footnotesize,at={(0.98,0.98)},anchor=north east}]
\addplot[blue,mark=*] table[x=N,y=L1_Bperp] {figures/data/alfven_conv.dat};
\addplot[black,densely dotted,domain=16:128,samples=2] {37.2*x^(-2)};
\legend{AMReX (MUSCL+SSP-RK2), наклон $-2$}
\end{loglogaxis}
\end{tikzpicture}
\caption{На гладком решении схема действительно второго порядка: наклон
графика $1.94$ против номинальных~2, тогда как первопорядковая схема ВКР на этой
же задаче даёт $1.2$--$1.3$. Именно здесь порядок и проверяется --- на разрывных
задачах он ограничен первым независимо от схемы. Циркулярно-поляризованная
альфвеновская волна, $t=1$ (полный период), CFL~0.4. Источник:
\code{docs/figures/data/alfven\_conv.dat}.}
\label{fig:conv}
\end{figure}

\paragraph{Задача Даи--Вудварда.}
Добавлена в набор в рамках настоящей работы (ВКР \S2.1.2): одномерная задача
Римана с двумя быстрыми и двумя медленными МГД-разрывами, двумя вращательными
и контактным. На $N=400$, $t=0.2$ решение конечно, положительно, монотонно;
диапазоны $\rho\in[1.00,1.66]$, $p\in[0.95,1.98]$, $B_y\in[1.0,1.6]$ совпадают
с рис.~7 в~\cite{shamanov2025}. Прямое сравнение с \code{legacy\_corrected} ---
в~\S\ref{sec:res-cmp}.

\paragraph{Вихрь Орзага--Танга и вращающийся цилиндр.}
На сетке $N=128$, $t_{\text{ОТ}}=0.5$, $t_{\text{цил}}=0.15$ диапазоны величин
совпадают с эталонными расчётами \cite{avdeeva2019,shamanov2025} с точностью
$\sim2\%$: для ОТ $\rho\in[0.091,0.492]$ (эталон $[0.087,0.489]$),
$p\in[0.028,0.505]$ (эталон $[0.028,0.49]$); для цилиндра
$\rho\in[0.66,11.55]$ (эталон $[0.571,10.60]$), $p\in[0.039,1.73]$. Срез
давления вдоль $y=0.3125$ (рис.~\ref{fig:ot}) воспроизводит структуру
рис.~8 в~\cite{avdeeva2019} по числу и положению признаков; паразитных
осцилляций не наблюдается. Диагональ $x=y$ для цилиндра
(рис.~\ref{fig:rotor}) симметрична, содержит два пика плотности и эвакуацию
давления в ядре --- как рис.~6 в~\cite{avdeeva2019} и рис.~20, 22
в~\cite{shamanov2025}; решатель первого порядка сильнее размывает пик
плотности.

\begin{figure}[htbp]\centering
\begin{minipage}{0.49\textwidth}\centering
\begin{tikzpicture}
\begin{axis}[width=\textwidth,height=4.6cm,xlabel=$x$,ylabel=$p$,xmin=0,xmax=1,
  title={\footnotesize ОТ, срез $y=0.3125$}]
\addplot[blue,mark=none] table[x=x,y=p] {figures/data/ot_slice.dat};
\end{axis}
\end{tikzpicture}
\caption{Срез давления в вихре Орзага--Танга воспроизводит структуру
рис.~8 в~\cite{avdeeva2019} по числу и положению особенностей; паразитных
осцилляций между ними нет. $N=128$, $t=0.5$, $y=0.3125$. Источник:
\code{docs/figures/data/ot\_slice.dat}.}
\label{fig:ot}
\end{minipage}\hfill
\begin{minipage}{0.49\textwidth}\centering
\begin{tikzpicture}
\begin{axis}[width=\textwidth,height=4.6cm,xlabel=$x$,ylabel=$\rho$,xmin=0,xmax=1,
  legend style={font=\tiny,at={(0.5,0.98)},anchor=north},
  title={\footnotesize цилиндр, диагональ $x{=}y$}]
\addplot[red,mark=none] table[x=x,y=rho] {figures/data/rotor_diag_amrex.dat};
\addplot[blue,densely dashed,mark=none] table[x=x,y=rho] {figures/data/rotor_diag_legacy.dat};
\legend{AMReX N=128,legacy\_corrected}
\end{axis}
\end{tikzpicture}
\caption{Вращающийся цилиндр: оба решателя дают симметричный профиль с двумя
пиками плотности и эвакуацией в ядре, но первопорядковый заметно ниже срезает
пики --- та же диссипация, что и в остальных тестах. $t=0.15$, диагональ $x=y$,
48~интервалов усреднения. Источник:
\code{docs/figures/data/rotor\_diag\_*.dat}.}
\label{fig:rotor}
\end{minipage}
\end{figure}

\begin{figure}[htbp]\centering
\begin{minipage}{0.49\textwidth}\centering
\begin{tikzpicture}
\begin{axis}[width=0.80\textwidth,axis equal image,enlargelimits=false,axis on top,
  xlabel=$x$,ylabel=$y$,title={\footnotesize Орзага--Танга, $\rho$, $t=0.5$},
  colormap/jet,colorbar,point meta min=0.0913,point meta max=0.4915,
  colorbar style={width=0.16cm,tick label style={font=\tiny}}]
\addplot graphics [xmin=0,xmax=1,ymin=0,ymax=1] {figures/maps/ot_rho.png};
\end{axis}
\end{tikzpicture}
\end{minipage}\hfill
\begin{minipage}{0.49\textwidth}\centering
\begin{tikzpicture}
\begin{axis}[width=0.80\textwidth,axis equal image,enlargelimits=false,axis on top,
  xlabel=$x$,ylabel=$y$,title={\footnotesize Орзага--Танга, $p$, $t=0.5$},
  colormap/jet,colorbar,point meta min=0.0281,point meta max=0.5045,
  colorbar style={width=0.16cm,tick label style={font=\tiny}}]
\addplot graphics [xmin=0,xmax=1,ymin=0,ymax=1] {figures/maps/ot_p.png};
\end{axis}
\end{tikzpicture}
\end{minipage}

\medskip
\begin{minipage}{0.49\textwidth}\centering
\begin{tikzpicture}
\begin{axis}[width=0.80\textwidth,axis equal image,enlargelimits=false,axis on top,
  xlabel=$x$,ylabel=$y$,title={\footnotesize Цилиндр, $\rho$, $t=0.15$},
  colormap/jet,colorbar,point meta min=0.5885,point meta max=11.554,
  colorbar style={width=0.16cm,tick label style={font=\tiny}}]
\addplot graphics [xmin=0,xmax=1,ymin=0,ymax=1] {figures/maps/rotor_rho.png};
\end{axis}
\end{tikzpicture}
\end{minipage}\hfill
\begin{minipage}{0.49\textwidth}\centering
\begin{tikzpicture}
\begin{axis}[width=0.80\textwidth,axis equal image,enlargelimits=false,axis on top,
  xlabel=$x$,ylabel=$y$,title={\footnotesize Цилиндр, $|\B|^2/2$, $t=0.15$},
  colormap/jet,colorbar,point meta min=0.0549,point meta max=2.509,
  colorbar style={width=0.16cm,tick label style={font=\tiny}}]
\addplot graphics [xmin=0,xmax=1,ymin=0,ymax=1] {figures/maps/rotor_pmag.png};
\end{axis}
\end{tikzpicture}
\end{minipage}

\medskip
\begin{minipage}{0.49\textwidth}\centering
\begin{tikzpicture}
\begin{axis}[width=0.72\textwidth,axis equal image,enlargelimits=false,axis on top,
  xlabel=$x$,ylabel=$y$,title={\footnotesize Петля поля, $|\B_\perp|$, $t=2$},
  colormap/jet,colorbar,point meta min=0,point meta max=0.001125,
  colorbar style={width=0.16cm,tick label style={font=\tiny},
    /pgf/number format/sci}]
\addplot graphics [xmin=-1,xmax=1,ymin=-0.5,ymax=0.5] {figures/maps/loop_bmag.png};
\end{axis}
\end{tikzpicture}
\end{minipage}\hfill
\begin{minipage}{0.49\textwidth}\centering
\begin{tikzpicture}
\begin{axis}[width=0.80\textwidth,axis equal image,enlargelimits=false,axis on top,
  xlabel=$x$,ylabel=$y$,title={\footnotesize CP-альфвен, $B_z$, $t=1$},
  colormap/jet,colorbar,point meta min=-0.0990,point meta max=0.0990,
  colorbar style={width=0.16cm,tick label style={font=\tiny}}]
\addplot graphics [xmin=0,xmax=1.1547,ymin=0,ymax=2] {figures/maps/alfven_bz.png};
\end{axis}
\end{tikzpicture}
\end{minipage}
\caption{Все четыре канонические задачи доходят до конца и дают структуру,
совпадающую с литературной: центральный ромб и диагональные фронты в вихре
Орзага--Танга, симметричные пики плотности у цилиндра, сохранившая форму петля
после двух проходов области, неискажённая альфвеновская волна после полного
периода. Диапазоны на шкалах совпадают с эталонными расчётами с точностью
$\sim2\,\%$ (\S\ref{sec:res-lit}). Однородная сетка, MUSCL + SSP-RK2 + CT;
Орзага--Танг и цилиндр --- $128^2$, петля --- $128\times64$, альфвеновская
волна --- $64\times111$. Источник: \code{scripts/field\_map.py}.}
\label{fig:maps}
\end{figure}


\subsection{Порядок сходимости интегратора и схемы}
\label{sec:res-conv}

Множитель перехода SSP-RK2 для линейного оператора совпадает с
$1+a+\tfrac12a^2$ с точностью $10^{-14}$; наблюдаемый порядок на скалярной ОДУ
$y'=\lambda y$ равен $2.004$\Pv. Наблюдаемый порядок циркулярно-поляризованной
альфвеновской волны --- $1.94$ по относительной $L_1$-ошибке $B_\perp$ и
$1.92$ по суммарной норме $L_1(\text{sum4})$, которую проверяет автотест
($N{=}64\to128$), при пороге приёмки $\geq1.8$\Iz.

\subsection{Консервативность на адаптивной сетке}
\label{sec:res-amr}

Иерархия уровней продвигается без подциклирования: шаг $\Delta t$ одинаков на
всех уровнях. Поэтому грубая ячейка, примыкающая к измельчённому блоку,
обновлялась \emph{грубым} потоком, тогда как мелкие ячейки за той же гранью ---
набором \emph{мелких} потоков. Эти два вклада ничем не согласовывались, и
иерархия не была консервативной. Введён регистр потоков
(\code{amrex::YAFluxRegister}) на каждую пару уровней: за стадию Эйлера
накапливаются грубый и мелкий вклады, а поправка
$\Delta U=\pm\Delta t/\Delta x_c\,(F_c-\langle F_f\rangle)$ применяется к грубым
ячейкам у стыка до \code{average\_down}. Компоненты $B_x, B_y$ из рефлюкса
\emph{исключены}: плоскостное поле ведёт constrained transport через узловые
ЭДС, и рефлюкс газового потока разрушил бы $\divg\B=0$.

Измерение на статической двухуровневой иерархии (вихрь Орзага--Танга, база
$64^2$, $r=2$, мелкий уровень покрывает половину базовой сетки, $t=0.05$, все
границы периодические --- через внешнюю границу не вытекает ничего, поэтому
$\sum\rho$ и $\sum e$ обязаны сохраняться точно):

\begin{center}
\begin{tabular}{lrr}
\toprule
Величина & без рефлюкса & с рефлюксом \\
\midrule
Дрейф $\sum\rho$   & $6.9\!\cdot\!10^{-5}$ & $6.3\!\cdot\!10^{-15}$ \\
Дрейф $\sum e$     & $1.5\!\cdot\!10^{-4}$ & $7.9\!\cdot\!10^{-16}$ \\
\bottomrule
\end{tabular}
\end{center}
\noindent\footnotesize
Источник: \code{tests/check\_amr\_conservation.py} (тест \code{amr.conservation}),
конфигурация \code{inputs/amr\_conservation.json}\Iz. Тест защищён от вырождения:
он падает, если мелкий уровень покрывает весь периодический домен (тогда стыка
уровней нет и рефлюкс тождественно ничего не делает), и повторяет прогон с
\code{reflux=false}, требуя, чтобы дефект действительно проявился.\normalsize

\medskip
При \emph{включённом} перестроении сетки (\code{regrid\_int}=4) дрейф равен
$9.7\!\cdot\!10^{-12}$ по массе и $7.0\!\cdot\!10^{-8}$ по энергии. Источник
остатка установлен измерением, а не предположением: само перестроение
консервативно --- скачок массы и полной энергии на одном \code{regrid} равен
$2.5\!\cdot\!10^{-16}$ и $3.2\!\cdot\!10^{-16}$ (тест
\code{amr.regrid\_conservation}). Расходится \emph{разделение} энергии:
полная энергия переносится консервативным скаляром, а поле --- независимым
бездивергентным оператором, поэтому магнитная часть $\tfrac12|\B|^2$ скачет на
$\sim\!10^{-5}$ относительных, и эту разницу поглощает тепловая часть, то есть
давление.

Проверен и обратный вариант (сохранять газовую энергию, флаг
\code{amr.regrid\_preserve\_pressure}): он \emph{не} уменьшает скачок тепловой
энергии ($1.58\!\cdot\!10^{-5}$ против $1.67\!\cdot\!10^{-5}$), поскольку
нелинейна и кинетическая часть $\tfrac12\rho|\vv|^2$, а полную энергию ухудшает
примерно в 25 раз. Принят вариант «сохраняем полную энергию» (решение~D-007).
Эффект убывает с разрешением: итоговый дрейф энергии $7.0\!\cdot\!10^{-8}$ при
$N=64$ против $1.5\!\cdot\!10^{-9}$ при $N=128$, то есть это ошибка
аппроксимации интерполяции, а не дефект консервативности.

\subsection{Независимый эталон Брио--Ву и пофронтовые метрики}
\label{sec:res-briowu-t06}

Эталон, посчитанный той же схемой, что и проверяемый результат, доказывает лишь
самосогласованность, поэтому прежнее сравнение было помечено как провизорное.
Точный решатель Римана здесь неприменим: решение содержит \emph{составную} волну
(медленная ударная с примыкающим разрежением), а классические семиволновые
решатели строятся в предположении регулярных волн.

Эталон построен схемой, не имеющей с проверяемой ни одной общей детали:
центральная схема Куртганова--Тадмора. Она не раскладывает решение по волнам и
вообще не решает задачу Римана --- нужен только физический поток и оценка
максимальной скорости; файл \code{tests/briowu\_reference.cpp} намеренно не
подключает \code{src/kernels}, чтобы ошибка в общем коде не воспроизвелась
одинаково в обеих ветвях. Эталон сходится (L1 относительно $N=6400$:
$2.68\!\cdot\!10^{-2}$, $1.27\!\cdot\!10^{-2}$, $5.41\!\cdot\!10^{-3}$,
$1.92\!\cdot\!10^{-3}$ при $N=400\ldots3200$; порядок $\approx1$, как и
ожидается на разрывном решении).

Взаимная проверка: KT при $N=6400$ против HLLD+MUSCL при $N=2048$ даёт
относительную $L_1$ $6.6\!\cdot\!10^{-4}$ ($\rho$), $1.7\!\cdot\!10^{-3}$ ($u$),
$6.4\!\cdot\!10^{-4}$ ($p$), $9.9\!\cdot\!10^{-4}$ ($B_y$); экстремумы совпадают
до четвёртой значащей цифры ($\rho_{\min}$ 0.11681 против 0.11692). Две схемы
без общего кода сходятся к одному слабому решению, включая составную волну\Iz.

\begin{center}\footnotesize\setlength{\tabcolsep}{5pt}
\begin{tabular}{lcccc}
\toprule
Схема ($N=400$) & $\rho$ & $u$ & $p$ & $B_y$ \\
\midrule
\code{legacy\_corrected}, HLLD, 1-й порядок & $1.19\!\cdot\!10^{-2}$ & $3.76\!\cdot\!10^{-2}$ & $1.28\!\cdot\!10^{-2}$ & $1.58\!\cdot\!10^{-2}$ \\
AMReX N0 (const + Эйлер + BS)               & $1.13\!\cdot\!10^{-2}$ & $3.36\!\cdot\!10^{-2}$ & $1.13\!\cdot\!10^{-2}$ & $1.40\!\cdot\!10^{-2}$ \\
AMReX N3 (MUSCL + SSP-RK2 + GS)             & $3.01\!\cdot\!10^{-3}$ & $1.07\!\cdot\!10^{-2}$ & $2.83\!\cdot\!10^{-3}$ & $4.52\!\cdot\!10^{-3}$ \\
\bottomrule
\end{tabular}
\end{center}
\noindent\footnotesize
Относительная $L_1$ относительно эталона KT, $N=6400$. Источник:
\code{docs/BRIOWU\_T06.md}, \code{benchmarks/summary/briowu\_t06\_*.json}\Iz.\normalsize

\medskip
Отсюда два вывода. Во-первых, \code{legacy\_corrected} и AMReX~N0 совпадают в
пределах 6--13\,\% --- теперь уже против независимого эталона, что подтверждает
эквивалентность реализаций HLLD. Во-вторых, выигрыш AMReX~N3 в 3--4 раза по
$L_1$ и в 15--30 раз по ширине фронта (1.5 ячейки против 46 на контактном
разрыве) даёт именно второй порядок, а не решатель Римана. Цена честная: у N3
появляются перелёты до $1.0\!\cdot\!10^{-2}$ по $u$ и больше смен знака первой
разности (26--33 против 2--6), которых у первопорядковых схем нет.

\subsection{Прямое сравнение схем на общих задачах}
\label{sec:res-cmp}

Сравнение охватывает подмножество тестового набора ВКР
\cite{shamanov2025}, которое обе программы решают на общей постановке:
одномерные задачи Римана Брио--Ву и Даи--Вудварда (гл.~2 ВКР), а также
двумерные задачи о вихре Орзага--Танга, вращающемся цилиндре и петле поля
(\S4.1 ВКР). Осесимметричные задачи (цилиндрическая задача Сода, МГД-взрыв в
$(r,z)$-геометрии) вне рамок сравнения: у \code{mhd2d\_amrex} нет
осесимметричного режима. Полный разбор ---
\code{docs/CROSS\_SOLVER\_COMPARISON.md}.

\paragraph{Одномерные задачи Римана: сопоставление реализаций HLLD.}
Драйвер \code{scripts/legacy\_1d\_riemann.cpp} собирается против исправленного
потока \code{HLLD\_flux\_corrected} на \emph{той же} равномерной одномерной
сетке, что и режим N0 нового решателя (кусочно-постоянно + Эйлер + HLLD).
Разность профилей --- это разность реализаций потока и сетки, но не порядка.
Эталон --- провизорный сошедшийся расчёт AMReX~N3, $N=2048$.

\begin{center}\footnotesize\setlength{\tabcolsep}{4.5pt}
\begin{tabular}{llcccc}
\toprule
Задача & Профиль & $L_1(\rho)$ & $L_1(B_y)$ & избыток TV$(B_y)$ & over/under \\
\midrule
\multirow{3}{*}{Брио--Ву}
 & \code{legacy\_corrected} HLLD & $1.14\!\cdot\!10^{-2}$ & $1.52\!\cdot\!10^{-2}$ & $0$ & $0/0$ \\
 & \code{mhd2d\_amrex} N0        & $1.08\!\cdot\!10^{-2}$ & $1.33\!\cdot\!10^{-2}$ & $0$ & $0/0$ \\
 & \code{mhd2d\_amrex} N3        & $2.52\!\cdot\!10^{-3}$ & $3.96\!\cdot\!10^{-3}$ & $0.11$ & мал.\ \\
\midrule
\multirow{3}{*}{Даи--Вудвард}
 & \code{legacy\_corrected} HLLD & $1.23\!\cdot\!10^{-2}$ & $1.43\!\cdot\!10^{-2}$ & $0$ & $0/0$ \\
 & \code{mhd2d\_amrex} N0        & $1.15\!\cdot\!10^{-2}$ & $1.33\!\cdot\!10^{-2}$ & $0$ & $0/0$ \\
 & \code{mhd2d\_amrex} N3        & $2.71\!\cdot\!10^{-3}$ & $3.15\!\cdot\!10^{-3}$ & $3.9\!\cdot\!10^{-3}$ & $0/0$ \\
\bottomrule
\end{tabular}
\end{center}
\noindent\footnotesize $N=400$, замороженные ГУ. На истинно одномерной сетке
исправленный поток \code{legacy\_corrected} и ядро \code{mhd2d\_amrex}
совпадают в пределах $\sim6\text{--}10\%$ во всех нормах на обеих задачах ---
реализации HLLD численно эквивалентны. Обе схемы N0 строго монотонны (нулевой
избыток вариации, отсутствие overshoot/undershoot); ни одна не даёт
осцилляций на вращательных разрывах, о которых ВКР сообщает для HLLC-L.
Выигрыш N3 над N0 ($\approx4.5\times$ по $L_1$) --- это второй порядок
пространства и времени, а не поток (рис.~\ref{fig:dw}). Источник:
\code{briowu\_1d\_flux\_comparison.json},
\code{dai\_woodward\_1d\_comparison.json}.\normalsize

\begin{figure}[htbp]\centering
\begin{tikzpicture}
\begin{axis}[width=0.49\textwidth,height=4.8cm,xlabel=$x$,ylabel=$\rho$,
  xmin=-0.5,xmax=0.5,legend style={font=\tiny,at={(0.02,0.98)},anchor=north west},
  title={\footnotesize Даи--Вудвард, плотность}]
\addplot[gray,thick] table[x=x,y=rho] {figures/data/dw1d_ref.dat};
\addplot[blue,densely dashed] table[x=x,y=rho] {figures/data/dw1d_legacy.dat};
\addplot[red] table[x=x,y=rho] {figures/data/dw1d_n3.dat};
\legend{эталон,legacy HLLD,AMReX N3}
\end{axis}
\begin{axis}[xshift=0.5\textwidth,width=0.49\textwidth,height=4.8cm,xlabel=$x$,
  ylabel=$p$,xmin=-0.5,xmax=0.5,title={\footnotesize давление}]
\addplot[gray,thick] table[x=x,y=p] {figures/data/dw1d_ref.dat};
\addplot[blue,densely dashed] table[x=x,y=p] {figures/data/dw1d_legacy.dat};
\addplot[red] table[x=x,y=p] {figures/data/dw1d_n3.dat};
\end{axis}
\end{tikzpicture}
\caption{Даи--Вудвард: задача с семью волнами, включая две вращательные,
--- проверка того, что выигрыш второго порядка не ограничен одной удобной
постановкой. При равной сетке $N=400$ ошибка падает в 4.5~раза, а все семь
разрывов остаются на своих местах; диапазоны $\rho\in[1.0,1.66]$,
$p\in[0.95,1.98]$ совпадают с рис.~7 в~\cite{shamanov2025}. $t=0.2$. Источник:
\code{docs/figures/data/dw1d\_*.dat}.}
\label{fig:dw}
\end{figure}

\paragraph{Аблация и стоимость (Брио--Ву, полоса).}
Отдельная серия на 2D-полосе разделяет вклад декартовой сетки, MUSCL и выбора
CT/интегратора и даёт диагностические времена:

\begin{center}\footnotesize\setlength{\tabcolsep}{5pt}
\begin{tabular}{lccc}
\toprule
Профиль & $L_1(\rho)$ & избыток TV$(B_y)$ & время, с \\
\midrule
\code{legacy\_corrected} $400\times8$ (полоса) & $4.3\!\cdot\!10^{-2}$ & $0.68$ & $209$ \\
N1 (MUSCL + Эйлер + BS)  & $2.4\!\cdot\!10^{-3}$ & $0.30$ & --- \\
N2 (MUSCL + SSP-RK2 + BS)& $2.6\!\cdot\!10^{-3}$ & $0.18$ & --- \\
N3 (MUSCL + SSP-RK2 + GS)& $2.5\!\cdot\!10^{-3}$ & $0.11$ & $0.94$ \\
\bottomrule
\end{tabular}
\end{center}
\noindent\footnotesize Большая ошибка \code{legacy\_corrected} на полосе (в
$\sim4$ раза относительно N0) --- это неструктурированная сетка (RT0, узловая
ЭДС, поперечный разброс $\sim20\%$), а не поток. SSP-RK2 и Gardiner--Stone
уменьшают избыток вариации ($0.30\to0.18\to0.11$) без потери точности. При
равном шаге $\Delta x=1/128$ новая схема $\approx113\times$ быстрее, при
$\Delta x=1/400$ --- $\approx222\times$ (диагностика, не benchmark). Источник:
\code{CROSS\_SOLVER\_BRIOWU\_1D.md}.\normalsize

\paragraph{Задача о петле магнитного поля.}
Область $[-1,1]\times[-0.5,0.5]$, $\rho=p=1$, $\vv=(2,1)$, $t=2$ (два прохода
домена). Метрика --- $E_B(t)/E_B(0)$ (единица для недиссипативной схемы). Это
именно тот тест, на котором ВКР \cite{shamanov2025} (\S4.1.4) обнаруживает
<<свойство высокой диссипации выбранного метода>> и рекомендует <<выбрать
численный метод более высокого порядка>>. Сравнение даёт количественную меру:

\begin{center}\footnotesize
\begin{tabular}{lccc}
\toprule
Решатель & Сетка & $E_B(t)/E_B(0)$ & нормир.\ $\|\divg\B\|$ \\
\midrule
\code{legacy\_corrected} (T05)      & $64\times32$   & $0.176$ & $2.2\!\cdot\!10^{-15}$ \\
\code{legacy\_corrected}            & $128\times64$  & $0.312$ & $4.0\!\cdot\!10^{-15}$ \\
\code{mhd2d\_amrex}                 & $128\times64$  & $\mathbf{0.836}$ & $7\!\cdot\!10^{-16}$ \\
\code{mhd2d\_amrex}                 & $256\times128$ & $0.913$ & --- \\
\bottomrule
\end{tabular}
\end{center}
\noindent\footnotesize При равной сетке $128\times64$ первопорядковый решатель
сохраняет $\approx31\%$ магнитной энергии петли за два прохода,
второпорядковый --- $\approx84\%$ (и $\to1$ с измельчением, рис.~\ref{fig:loop});
рост $\divg\B$ у обоих --- до уровня округления, что совпадает с
\S4.1.4 в~\cite{shamanov2025} ($5.9\!\cdot\!10^{-17}\to2.0\!\cdot\!10^{-15}$).
\normalsize

\begin{figure}[htbp]\centering
\begin{tikzpicture}
\begin{axis}[width=0.62\textwidth,height=4.6cm,xlabel={$N$ (по длинной стороне)},
  ylabel={$E_B(t{=}2)/E_B(0)$},ymin=0,ymax=1.05,
  xtick={64,128,256},xticklabels={64,128,256},
  legend style={font=\footnotesize,at={(0.98,0.4)},anchor=east}]
\addplot[red,mark=*] table[x=N,y=eb_ratio] {figures/data/loop_eb_amrex.dat};
\addplot[blue,mark=square*,only marks] coordinates {(64,0.176) (128,0.312)};
\legend{\code{mhd2d\_amrex}, \code{legacy\_corrected}}
\end{axis}
\end{tikzpicture}
\caption{Тест, на котором ВКР констатирует «высокую диссипацию» метода, и
мера этой диссипации. Идеальная схема сохранила бы единицу. При равной сетке
$128\times64$ первопорядковая схема удерживает $31\,\%$ магнитной энергии
петли за два прохода области, второпорядковая --- $84\,\%$, и с измельчением
приближается к единице. Источник:
\code{docs/figures/data/loop\_eb\_amrex.dat}.}
\label{fig:loop}
\end{figure}

\paragraph{Вихрь Орзага--Танга.} Прямое сравнение при равном разрешении
затруднено: \code{legacy\_corrected} на структурной сетке $N=128$ не доходит
до $t=0.5$ (\S\ref{sec:res-legacy}); в \cite{shamanov2025,avdeeva2019}
использованы сетки $N\approx400\text{--}800$ на неправильной триангуляции.
Расчёт \code{legacy\_corrected} на $N=256$ (число Куранта $0.2$;
$\sim1.3\!\cdot\!10^5$ треугольников) выполняется --- это следующий шаг
верификации legacy. Диапазоны и срез \code{mhd2d\_amrex} совпадают с
литературой (\S\ref{sec:res-lit}).

\subsection{Происхождение перелётов у фронтов и выбор лимитера}
\label{sec:res-rp1}

У схемы второго порядка на Брио--Ву видны перелёты. Чтобы указать их источник,
собрана матрица прогонов, в которой за раз меняется ровно одна деталь схемы;
все варианты оцениваются относительно независимого эталона
(\S\ref{sec:res-briowu-t06}), $N_x=400$, $t=0.1$, CFL 0.1. Ширина фронта и
амплитуда выброса приводятся раздельно намеренно: лимитер, снимающий
осцилляции ценой вдвое размазанного контактного разрыва, улучшением не
является.

\begin{center}\footnotesize\setlength{\tabcolsep}{5pt}
\begin{tabular}{lcccc}
\toprule
Вариант & отн. $L_1$ & перелёт & смен знака & контакт, яч. \\
\midrule
N0 кусочно-пост. + Эйлер + BS & $1.75\!\cdot\!10^{-2}$ & $4.4\!\cdot\!10^{-3}$ & 18 & 14.24 \\
N1 {}+ MUSCL-MC               & $5.75\!\cdot\!10^{-3}$ & $\mathbf{3.26\!\cdot\!10^{-2}}$ & 181 & 1.66 \\
N2 {}+ SSP-RK2                & $5.33\!\cdot\!10^{-3}$ & $1.84\!\cdot\!10^{-2}$ & 154 & 1.85 \\
N3 {}+ Gardiner--Stone        & $5.26\!\cdot\!10^{-3}$ & $1.18\!\cdot\!10^{-2}$ & 118 & 1.49 \\
\midrule
minmod                        & $8.02\!\cdot\!10^{-3}$ & $7.6\!\cdot\!10^{-3}$ & 80 & 2.53 \\
van Leer                      & $5.91\!\cdot\!10^{-3}$ & $9.4\!\cdot\!10^{-3}$ & 114 & 2.01 \\
MC (рабочий)                  & $5.26\!\cdot\!10^{-3}$ & $1.18\!\cdot\!10^{-2}$ & 118 & 1.49 \\
\bottomrule
\end{tabular}
\end{center}
\noindent\footnotesize
Источник: \code{docs/RP1\_MONOTONICITY.md},
\code{benchmarks/summary/rp1\_limiter\_study\_n400.json}\Iz.\normalsize

\medskip
Перелёт вносит \emph{реконструкция}: переход N0\,$\to$\,N1 меняет только её и
увеличивает выброс в 7.4 раза, одновременно втрое уменьшая ошибку и сужая
контактный разрыв с 14.2 до 1.7 ячейки. Интегратор и усреднение ЭДС выброс не
добавляют, а частично снимают (до $1.18\!\cdot\!10^{-2}$), поэтому итоговая
схема монотоннее, чем MUSCL сам по себе.

Выбор лимитера --- измеренный размен, а не предпочтение: minmod даёт выброс в
1.55 раза меньше, но проигрывает MC 1.52 раза по точности и в 1.7 раза по
резкости контакта. Ни один из трёх не выигрывает по обоим показателям.

Отдельно проверена гипотеза, что перелёт по скорости порождается ограничением
примитивных $u,p$ вместо консервативных $\rho u, e$. Она \textbf{не
подтвердилась}: реконструкция по консервативным переменным меняет выброс на
$4\,\%$ ($1.182\!\cdot\!10^{-2}\to1.137\!\cdot\!10^{-2}$), ухудшая точность на
$47\,\%$ и расширяя контакт в 1.49 раза, причём на 120 гранях восстановленное
состояние оказывается недопустимым и грань считается по первому порядку.
Рабочим остаётся примитивный набор\Iz.

\medskip
Выбор SSP-RK2 проверен отдельно, потому что утверждение о нём легко сделать
сильнее, чем позволяет измерение. SSP-свойство --- это утверждение о том, до
какого шага по времени TVD-оценка шага Эйлера переносится на шаг целиком, и
проверять его при CFL~0.1 бессмысленно: там обе схемы заведомо далеко внутри
области, где разницы не ожидается. Развёртка по всему допустимому шагу:

\begin{center}\footnotesize\setlength{\tabcolsep}{6pt}
\begin{tabular}{ccccc}
\toprule
& \multicolumn{2}{c}{SSP-RK2 (Хойн)} & \multicolumn{2}{c}{RK2 средней точки (не SSP)} \\
\cmidrule(lr){2-3}\cmidrule(lr){4-5}
CFL & отн.\ $L_1$ & перелёт & отн.\ $L_1$ & перелёт \\
\midrule
0.1 & $5.264\!\cdot\!10^{-3}$ & $1.182\!\cdot\!10^{-2}$ & $5.264\!\cdot\!10^{-3}$ & $1.183\!\cdot\!10^{-2}$ \\
0.4 & $5.311\!\cdot\!10^{-3}$ & $1.154\!\cdot\!10^{-2}$ & $5.281\!\cdot\!10^{-3}$ & $1.142\!\cdot\!10^{-2}$ \\
0.6 & $5.334\!\cdot\!10^{-3}$ & $1.177\!\cdot\!10^{-2}$ & $5.378\!\cdot\!10^{-3}$ & $1.186\!\cdot\!10^{-2}$ \\
0.8 & $5.388\!\cdot\!10^{-3}$ & $1.149\!\cdot\!10^{-2}$ & $5.385\!\cdot\!10^{-3}$ & $1.171\!\cdot\!10^{-2}$ \\
0.9 & $5.430\!\cdot\!10^{-3}$ & $1.127\!\cdot\!10^{-2}$ & $5.404\!\cdot\!10^{-3}$ & $1.161\!\cdot\!10^{-2}$ \\
\bottomrule
\end{tabular}
\end{center}
\noindent\footnotesize
Источник: \code{benchmarks/summary/rp2\_integrator\_cfl\_sweep.json}\Iz.\normalsize

\medskip
Ни один прогон не отказал, и разница между SSP и не-SSP вариантом нигде не
превышает $3\,\%$ по выбросу --- меньше, чем разброс между соседними значениями
CFL у одного и того же интегратора. Отрицательный результат, таким образом,
получен не в одной удобной точке, а на всём диапазоне шага. Правдоподобное
объяснение --- обе схемы двухстадийные с одинаковой первой стадией, а на разрыве
ограничитель и диссипация HLLD срезают возмущение раньше, чем разница весов
второй стадии успевает накопиться; доказать это измерения не позволяют, поэтому
объяснение остаётся объяснением.

% ---------------------------------------------------------------------
\subsection{Сравнение схем по общей метрике}
\label{sec:res-rp3}

Предыдущие сравнения велись по разным величинам для разных задач. Здесь обе
схемы сведены на одну метрику при равном разрешении.

ВКР приводит для Брио--Ву таблицу «усреднённой ошибки» на сетках $N=64\ldots512$,
но её определение в тексте не выписано. Поэтому оно зафиксировано явно и
применяется к обеим схемам одинаково:
\[
 \delta_N=\frac1N\sum_i\ \sum_{v\in\{\rho,u,p,B_y\}} \bigl|q_v(x_i)-q_v^{\text{эт}}(x_i)\bigr| .
\]
Эталон --- независимая схема Куртганова--Тадмора при $N=6400$, а не собственный
расчёт на мелкой сетке; почему это важно, сказано ниже.

\begin{center}\footnotesize\setlength{\tabcolsep}{5pt}
\begin{tabular}{rcccccc}
\toprule
$N$ & legacy (треуг.) & AMReX N0 (декарт.) & AMReX N3 & legacy/N3 & legacy/N0 & N0/N3 \\
\midrule
64  & 0.6251 & 0.1766 & 0.0749 & 8.3 & 3.5 & 2.4 \\
128 & 0.4353 & 0.1193 & 0.0443 & 9.8 & 3.6 & 2.7 \\
256 & 0.3918 & 0.0718 & 0.0206 & 19.0 & 5.5 & 3.5 \\
400 & 0.2911 & 0.0580 & 0.0172 & 16.9 & 5.0 & 3.4 \\
512 & 0.2516 & 0.0500 & 0.0140 & \textbf{18.0} & 5.0 & 3.6 \\
\bottomrule
\end{tabular}
\end{center}
\noindent\footnotesize
Источник: \code{docs/RP3\_SCHEME\_COMPARISON.md},
\code{benchmarks/summary/rp3\_delta\_n.json}\Iz.\normalsize

\medskip
Выигрыш раскладывается на две независимые части, и это главное, что даёт
таблица. Переход с треугольной сетки на декартову при \emph{неизменном} первом
порядке даёт множитель 5.0, а второй порядок --- ещё 3.4--3.6; вместе 17--19 раз.
Разложение возможно потому, что N0 --- буквально та же схема, что и
\code{legacy\_corrected}, только на другой сетке.

Множитель 5.0 мог оказаться артефактом: профиль legacy снимается с треугольной
полосы усреднением по $y$, и в $\delta_N$ попадает поперечный разброс, которого
у одномерного декартова расчёта нет. Проверено прямо --- проекция сохраняет
среднеквадратичное отклонение внутри каждого столбца, и его сумма по четырём
переменным составляет $7\text{--}8\,\%$ от $\delta_N$ на всех сетках. Основную
часть ошибки поперечный разброс не объясняет; локально, на фронтах, он всё же
велик --- до $21\,\%$ диапазона по скорости\Iz.

$\delta_N$ усредняется по области, поэтому схема может выиграть в ней на гладких
участках, почти не тронув разрывы. На вопрос «во сколько раз резче» отвечает
только ширина фронта. Для контактного разрыва ($x\approx0.56$, единственного
хорошо отделённого от соседей) при $N=400$: \code{legacy\_corrected} ---
36.8~ячейки, AMReX~N0 --- 14.1, AMReX~N3 --- \textbf{3.9}; эталон показывает, что
истинная ширина здесь меньше ячейки. \textbf{При равном разрешении новая схема
разрешает контактный разрыв в 9.4 раза резче старой}\Iz.

\paragraph{Почему эталон независимый, а не собственный.} Протокол ВКР ---
сравнение с собственным расчётом на $N=1024$ --- проверяет самосогласованность:
систематическая ошибка, одинаковая на всех сетках, из такой метрики выпадает.
Обе версии метрики посчитаны для новой схемы, и расхождение видно прямо:
наблюдаемый порядок при $N=512$ равен $0.83$ против независимого эталона и
$1.50$ против собственного. Чем ближе сетка к эталонной, тем сильнее завышение.
Это же объясняет, почему опубликованная в ВКР таблица показывает рост
наблюдаемого порядка с $0.49$ до $1.27$ при измельчении.

Наблюдаемый порядок около единицы у схемы второго порядка --- не дефект: на
разрывном решении $L_1$-порядок ограничен первым для любой схемы. Второй порядок
проверяется на гладком тесте (рис.~\ref{fig:conv}).

\subsection{Стоимость расчёта}
\label{sec:res-perf}

Замеры выполнены на одной рабочей станции без пиннинга и без эксклюзивного
доступа, поэтому они \emph{диагностические}: годятся для сравнения вариантов
между собой и не заменяют кластерную кампанию. Каждая запись содержит машинный
манифест и git-commit, один отбрасываемый прогрев и пять зачётных повторов;
отчитываются медиана и MAD. Все строки ниже сняты в одном режиме --- без записи
файлов результатов и в один поток, --- иначе нормированная стоимость между
случаями несравнима.

\begin{center}\footnotesize\setlength{\tabcolsep}{4pt}
\begin{tabular}{lccrrrr}
\toprule
Случай & сетка & ур. & шагов & медиана, с & мс/шаг & мкс/(яч$\cdot$шаг) \\
\midrule
Брио--Ву              & $512\times4$   & 0 & 391 & 0.5362 & 1.371 & 0.670 \\
CP-альфвен            & $64\times110$  & 0 & 128 & 0.4799 & 3.749 & 0.533 \\
Петля поля            & $128\times64$  & 0 & 527 & 1.4104 & 2.676 & \textbf{0.327} \\
Орзага--Танг          & $128^2$        & 0 & 323 & 2.2351 & 6.920 & 0.422 \\
Вращающийся цилиндр   & $128^2$        & 0 & 140 & 1.1876 & 8.483 & 0.518 \\
МГД-взрыв             & $128^2$        & 0 & 241 & 2.2189 & 9.207 & 0.562 \\
\midrule
ОТ + 1 уровень, стат. & $64^2$         & 1 & 37  & 0.2896 & 7.828 & 1.911 \\
ОТ + 1 уровень, regrid& $64^2$         & 1 & 37  & 0.3266 & 8.826 & 2.155 \\
Цилиндр + 1 уровень   & $200^2$        & 1 & 470 & 19.5675 & 41.633 & 1.041 \\
\bottomrule
\end{tabular}
\end{center}
\noindent\footnotesize
Разброс по всем строкам 0.6--1.7\,\%. Источник:
\code{docs/RP4\_RP5\_PERFORMANCE.md}, \code{benchmarks/summary/rp4\_timing\_*.json}\Iz.
\normalsize

\medskip
На однородной сетке стоимость ячейко-шага укладывается в 0.33--0.67~мкс при
разбросе задач от одномерной полосы до МГД-взрыва: патологически дорогих режимов
у схемы нет, и HLLD с откатом стоит примерно одинаково на гладком решении и на
сильном разрыве. Самая дорогая строка --- Брио--Ву, и дело не в физике, а в форме
области: полоса $512\times4$ при двух слоях фантомных ячеек тратит на границы
столько же, сколько на счёт внутри.

Адаптивная сетка обходится втрое-впятеро дороже в пересчёте на базовую ячейку,
но это не work-normalized величина --- мелкие уровни добавляют ячейки, которых
нет в знаменателе. Сравнивать между собой можно только AMR-строки: перестроение
сетки каждые четыре шага добавляет 13\,\%, а с ростом базовой сетки удельная
стоимость падает (1.91~мкс при базе $64^2$ против 1.04 при $200^2$) --- накладные
расходы на стык уровней амортизируются.

Профиль горячего пути (сборка с \code{AMReX\_TINY\_PROFILE}, вихрь
Орзага--Танга $64^2$ + один уровень AMR) объясняет, откуда эти накладные расходы
берутся. Inclusive-доли: стадия Эйлера 89.6\,\%, заполнение фантомов граневого
поля \textbf{29.7\,\%}, вычисление потоков и ЭДС 16.2\,\%, фантомы клеточных
величин 12.3\,\%, огрубление 8.5\,\%, согласование потоков 6.6\,\%, применение
обновлений 3.0\,\%\Iz. На небольшой AMR-задаче заполнение фантомов граневого поля
дороже самой физики --- почти вдвое. Согласование потоков стоит около 7\,\%:
недорого за восстановленную консервативность (\S\ref{sec:res-amr}).

\subsection{Параллельная эффективность на одном узле}
\label{sec:res-scaling}

Процессор рабочей станции неоднороден: 4 производительных ядра и 6
энергоэффективных. Без этого таблицы читаются неверно.

\begin{center}\footnotesize
\begin{tabular}{rccccrcc}
\toprule
& \multicolumn{3}{c}{OpenMP} & & \multicolumn{3}{c}{MPI} \\
\cmidrule(lr){2-4}\cmidrule(lr){6-8}
$p$ & медиана, с & $S(p)$ & $E(p)$ & & медиана, с & $S(p)$ & $E(p)$ \\
\midrule
1  & 2.2966 & 1.00 & 100\,\% & & 2.2871 & 1.00 & 100\,\% \\
2  & 1.3296 & 1.73 & 86\,\%  & & 1.2340 & 1.85 & 93\,\% \\
4  & 0.8530 & \textbf{2.69} & 67\,\% & & 0.6838 & \textbf{3.34} & \textbf{84\,\%} \\
6  & 0.9137 & 2.51 & 42\,\%  & & --- & --- & --- \\
8  & 0.9813 & 2.34 & 29\,\%  & & --- & --- & --- \\
10 & 1.0400 & 2.21 & 22\,\%  & & --- & --- & --- \\
\bottomrule
\end{tabular}
\end{center}
\noindent\footnotesize
Орзага--Танг $128^2$, однородная сетка, $t=0.5$, без записи файлов, 5 повторов.
Каждая конфигурация проверяется на совпадение итоговых диапазонов $\rho$ и $p$
с серийным прогоном: ускорение, изменившее ответ, ускорением не считается.
Источник: \code{benchmarks/summary/scaling\_\{omp,mpi\}\_ot128.json}\Iz.\normalsize

\medskip
Ускорение OpenMP достигает максимума ровно на четырёх потоках --- числе
производительных ядер; дальше работа уходит на энергоэффективные, и добавление
потоков замедляет расчёт. MPI на тех же четырёх ядрах эффективнее (3.34 против
2.69), и профиль объясняет почему: заполнение фантомов граневого поля занимает
почти треть шага, при MPI делится по рангам, а внутри ранга не распараллелено
вовсе. Независимость результата от разбиения по рангам закреплена автотестом\Pv.

Слабое масштабирование отвечает на другой вопрос --- не «во сколько раз
быстрее», а «сколько стоит рост задачи». Число ячеек на ранг (16384) и число
шагов (50) фиксированы, домен растёт вместе с числом рангов; идеал --- постоянное
время.

\begin{center}\footnotesize
\begin{tabular}{rccc}
\toprule
рангов & сетка & медиана, с & $E_{\text{weak}}$ \\
\midrule
1 & $128\times128$ & 0.3751 & 100\,\% \\
2 & $256\times128$ & 0.3876 & 97\,\% \\
4 & $256\times256$ & 0.4162 & \textbf{90\,\%} \\
8 & $512\times256$ & 0.8012 & 47\,\% \\
\bottomrule
\end{tabular}
\end{center}
\noindent\footnotesize
Источник: \code{benchmarks/summary/rp5\_weak\_ot.json}\Iz.\normalsize

\medskip
До четырёх рангов рост задачи почти бесплатен: обмены на границах доменов не
съедают выигрыш. Провал до 47\,\% на восьми рангах --- то же исчерпание
производительных ядер, что и в таблице OpenMP, а не свойство решателя. За
пределы одного узла слабое масштабирование на этой машине не измеримо в
принципе, и четыре точки выше --- отлаженный протокол, а не результат кампании.
Шаблоны заданий для очереди используют тот же измерительный скрипт, так что
кластерные числа будут сравнимы с приведёнными.

\subsection{Готовность к переносу на GPU}
\label{sec:res-gpu}

Конфигурация CUDA-сборки доходит до включения языка CUDA внутри AMReX и
останавливается на поиске \code{nvcc}: на доступной машине нет ни NVIDIA GPU, ни
CUDA-тулкита. Блокер внешний, поэтому вместо сборки выполнен аудит
GPU-безопасности\Iz.

Слой ядер к переносу готов: заголовки помечены device-квалификаторами, не
зависят от контейнеров AMReX (это гейтится автотестом), не выделяют память и не
используют виртуальных вызовов. Отдельно проверено самое подозрительное место ---
узловое усреднение ЭДС, на которое обычно и приходятся гонки. Оно безопасно:
цикл читает потоки на четырёх примыкающих гранях и пишет ровно одно значение на
узел, накопления нет, значит не нужны и атомарные операции. Разбиение по тайлам
не даёт двум тайлам владеть одним узлом, а на GPU тайлинг отключается.

Переделать придётся другое: 21 хостовый цикл заменяется на \code{ParallelFor};
редукции (шаг по времени, максимум $\divg\B$, интегралы диагностики) сейчас
захватывают host-переменную по ссылке и требуют перехода на \code{ReduceOps};
счётчики полов и откатов HLLD инкрементируются неатомарно. Последнее --- не
косметика: эти счётчики гейтят автотесты, и GPU-сборка не должна тихо потерять
диагностику, которой доверяет CPU-версия. Непроверяемый device-код в
репозиторий не добавлялся: на CPU он выродился бы в обычный цикл, прошёл бы
тесты и создал ложное впечатление готовности. Полный разбор и план первого
захода --- \code{docs/RP7\_GPU\_READINESS.md}.

\subsection{Сопоставление с литературой}
\label{sec:res-lit}

\begin{center}\footnotesize
\begin{tabular}{p{3.4cm}p{5.6cm}p{5.6cm}}
\toprule
Признак & Настоящая работа & Литература \\
\midrule
Брио--Ву, структура волн & составная (нерегулярная) ветвь: смена знака $B_y$,
плато $\rho\approx0.5$, контакт, медленный скачок
(\code{legacy\_corrected} и AMReX) & рис.~6 в~\cite{shamanov2025};
рис.~4 в~\cite{avdeeva2019} \\
CP-альфвен, порядок (1-й порядок) & \code{legacy\_corrected}:
$\approx1.2\text{--}1.3$ & табл.~3 в~\cite{shamanov2025}:
$R\approx1.22\text{--}1.29$ \\
Петля поля, рост $\divg\B$ & $1.4\!\cdot\!10^{-17}\to4.8\!\cdot\!10^{-16}$
(нормир.\ $2.2\!\cdot\!10^{-15}$) & \S4.1.4 в~\cite{shamanov2025}:
$5.9\!\cdot\!10^{-17}\to2.0\!\cdot\!10^{-15}$ \\
Вихрь ОТ, диапазоны (AMReX) & $\rho\in[0.091,0.492]$, $p\in[0.028,0.505]$
& рис.~7 в~\cite{avdeeva2019}: $\rho\in[0.087,0.489]$; рис.~25:
$p\in[0.028,0.49]$ \\
Цилиндр, диапазоны (AMReX) & $\rho\in[0.66,11.55]$, $p\in[0.039,1.73]$ &
рис.~5 в~\cite{avdeeva2019}: $\rho\in[0.571,10.60]$; рис.~19
в~\cite{shamanov2025}: $\rho\in[0.55,11]$ \\
CP-альфвен, порядок (2-й порядок) & AMReX: $1.94$ ($B_\perp$) / $1.92$ (sum4) & --- (в литературе только
1-й порядок) \\
\bottomrule
\end{tabular}
\end{center}

Не воспроизведены (вне рамок работы либо следующие фазы): таблицы усреднённых
ошибок ВКР 1--4 (показана согласованность порядка и величины, а не побитовое
воспроизведение); таблица времён CPU/GPU ВКР (иное оборудование, GPU-путь
отключён); независимое точное решение составной ветви Брио--Ву.

% =====================================================================
\section{Ограничения}
\label{sec:limits}

\begin{enumerate}[nosep]
\item \textbf{Давление слегка меняется при перестроении сетки.} Полная энергия
  переносится консервативным интерполятором, а поле --- независимым
  бездивергентным оператором; их магнитные части расходятся, и разницу поглощает
  тепловая энергия. Сдвиг составляет $\sim10^{-5}$ относительных на одно
  перестроение, убывает с разрешением и не устраняется сохранением давления
  вместо энергии --- это проверено, и обратный вариант оказался хуже вчетверо
  по полной энергии (\S\ref{sec:res-amr}).
\item \textbf{GPU-версия не собрана.} Не по внутренней причине: на доступной
  машине нет ни NVIDIA GPU, ни CUDA-тулкита, и конфигурация останавливается на
  поиске \code{nvcc}. Аудит переносимости выполнен, план первого захода
  составлен (\S\ref{sec:res-gpu}), но кода, который нельзя запустить и сверить,
  в репозитории нет.
\item \textbf{Многоузловое масштабирование не измерено.} На одном узле измерены
  и сильное, и слабое (\S\ref{sec:res-scaling}); кластерная кампания требует
  доступа, которого нет. Приведённые времена диагностические: одна рабочая
  станция, без пиннинга и эксклюзивного доступа.
\item \textbf{Эталон Брио--Ву --- численный, хотя и независимый.} Точное решение
  недостижимо: конфигурация содержит составную волну, которую классические
  семиволновые решатели не представляют. Эталон построен центральной схемой без
  решателя Римана, сошёлся и согласуется с проверяемой схемой до четвёртой
  значащей цифры по экстремумам (\S\ref{sec:res-briowu-t06}) --- но это всё же
  численное решение, а не аналитическое.
\item \textbf{Множитель 5.0 за декартову сетку получен на полосе.} Профиль
  старой схемы снимается с треугольной полосы усреднением по поперечной
  координате. Вклад поперечного разброса измерен и составляет $7\text{--}8\,\%$
  ошибки (\S\ref{sec:res-rp3}), так что вывод устоял, но полностью развести
  «сетку» и «усреднение» на этой постановке нельзя.
\item \textbf{Точность старой схемы на вихре Орзага--Танга ограничена первым
  порядком.} Расчёт доводится до конца, но при $N=128$ диффузия сужает
  диапазоны на 10--18\,\% относительно литературных при $N=800$. Это
  ограничение разрешения и порядка, а не устойчивости.
\item \textbf{Осесимметричные задачи выпускной работы} (цилиндрическая задача
  Сода, МГД-взрыв в $(r,z)$) в сравнение не входят: у нового решателя нет
  осесимметричного режима.
\item \textbf{Числа выпускной работы не считаются воспроизведёнными.} Показано
  согласие по порядку и диапазону, а не побитовое совпадение; часть
  расхождений объясняется тем, что там использовался self-reference
  (\S\ref{sec:res-rp3}).
\end{enumerate}

% =====================================================================
\section{Что осталось сделать}
\label{sec:plan}

Работы разделены по тому, что их блокирует.

\textbf{Ждут внешнего доступа.} Кампания на кластере: сильное и слабое
масштабирование за пределами одного узла. Шаблоны заданий готовы и вызывают тот
же измерительный скрипт, что и замеры здесь, поэтому кластерные числа будут
сравнимы с приведёнными без пересчёта. Перенос на GPU: замена хостовых циклов,
редукции через штатные механизмы, решение по счётчикам диагностики, затем
проверка совпадения CPU- и GPU-результата и бенчмарк. Порядок важен --- без
доступа обе задачи не начинаются, и именно его надо решать первым.

\textbf{Не заблокировано ничем.} Осесимметричный режим у нового решателя ---
он замкнул бы сравнение на полном тестовом наборе выпускной работы. Подциклирование
по времени на уровнях иерархии: сейчас все уровни идут общим шагом, и мелкий
уровень задаёт шаг всей сетке. Устранение сдвига давления при перестроении ---
для этого поле и энергию нужно интерполировать согласованно, а не независимо.

\textbf{Задел на будущее.} Границы между слоями «физика --- пространственный
оператор --- интегратор --- сетка --- диагностика» уже зафиксированы, и слой ядер
проверяется автотестом на независимость от библиотеки сетки. Это сделано,
чтобы разрывный метод Галёркина можно было добавить, не переписывая
конечно-объёмные регрессии.

% =====================================================================
\section{Заключение}

\textbf{Что построено.} Исторический решатель восстановлен и исправлен:
создан профиль с реестром из четырнадцати численных дельт, каждая с
регрессионным тестом. Разработан решатель второго порядка на декартовой
блочно-структурированной сетке с адаптивным измельчением --- MUSCL, SSP-RK2,
constrained transport, HLLD с положительность-сохраняющим откатом. Построен
воспроизводимый контур проверки достоверности: 21 автоматический тест,
манифесты запусков, машинно вычисляемые метрики.

\textbf{Что измерено.} Все канонические задачи исправленного решателя доходят до
конца, физически допустимы, консервативны и дискретно бездивергентны
(нормированная норма $\divg\B\leq6\!\cdot\!10^{-15}$); исходная ревизия на них
расходится. Новый решатель показывает порядок $1.94$ на гладком тесте,
воспроизводит составную ветвь задачи Брио--Ву и совпадает с эталонными
расчётами по диапазонам и структуре волн с точностью $\sim2\,\%$. На адаптивной
иерархии согласование потоков на стыке уровней восстанавливает консервативность:
дрейф массы падает с $6.9\!\cdot\!10^{-5}$ до уровня округления.

\textbf{Сравнение схем.} При равном разрешении на задаче Брио--Ву новая схема
даёт ошибку в 17--19 раз меньшую и разрешает контактный разрыв в 9.4 раза резче
(3.9 ячейки против 36.8). Выигрыш разложен на части: 5.0 раза даёт переход на
декартову сетку при неизменном первом порядке, ещё 3.4--3.6 --- второй порядок.
Отдельно проверено, что реализации HLLD в обоих решателях численно эквивалентны:
на истинно одномерной сетке они совпадают в пределах $6\text{--}10\,\%$, то есть
выигрыш даёт дискретизация, а не решатель Римана. На тесте с петлёй поля, где
выпускная работа и констатирует высокую диссипацию, новая схема сохраняет
$84\,\%$ магнитной энергии против $31\,\%$ при равной сетке.

\textbf{Отрицательные результаты, которые стоит назвать.} SSP-свойство
интегратора на задаче Брио--Ву не даёт ничего измеримого: SSP-RK2 и не-SSP схема
средней точки совпадают по ошибке до четвёртой значащей цифры на всём допустимом
шаге вплоть до $\mathrm{CFL}=0.9$. Гипотеза о том, что перелёты по скорости
порождаются реконструкцией примитивных переменных, не подтвердилась: переход к
консервативным меняет выброс на $4\,\%$, ухудшая точность на $47\,\%$.
Предположение, что преимущество декартовой сетки --- артефакт усреднения
треугольной полосы по поперечной координате, тоже не подтвердилось: вклад
разброса составляет $7\text{--}8\,\%$.

\textbf{Гипотеза H1 поддержана на однородной сетке.} Новая схема при равном шаге
устойчивее --- доводит до конца вихрь Орзага--Танга там, где первопорядковая
требует более подробной сетки, --- и на порядок точнее. Итоговым результатом это
не является: адаптивная сетка в сравнении схем не задействована, времена
диагностические, осесимметричные задачи вне сравнения.

\textbf{Не заявляется.} Готовность GPU-версии: сборка невозможна на доступном
оборудовании, выполнен аудит переносимости и составлен план. Многоузловое
масштабирование: на одном узле измерены сильное и слабое, кластерная кампания
требует доступа. Такой статус не прячет отрицательных и неполных результатов и
делает каждую следующую проверку воспроизводимой.

% =====================================================================
\begin{thebibliography}{99}
\bibitem{brio1988} Brio~M., Wu~C.\,C. An upwind differencing scheme for the
equations of ideal magnetohydrodynamics // J. Comput. Phys. 1988. Vol.~75.
P.~400--422.
\bibitem{miyoshi2005} Miyoshi~T., Kusano~K. A multi-state HLL approximate
Riemann solver for ideal magnetohydrodynamics // J. Comput. Phys. 2005.
Vol.~208. P.~315--344.
\bibitem{galanin2007} Галанин~М.\,П., Лукин~В.\,В. Разностная схема для решения
двумерных задач идеальной МГД на неструктурированных сетках. Препринт ИПМ
им.~М.\,В.~Келдыша РАН. № 50. 2007. 29~с.
\bibitem{avdeeva2019} Avdeeva~E.\,N., Lukin~V.\,V. Divergence-free
finite-difference method for 2D ideal MHD // J. Phys.: Conf. Ser. 2019.
Vol.~1336. 012026.
\bibitem{shamanov2025} Шаманов~И.\,П. Разработка и исследование программного
комплекса для решения задач идеальной магнитной гидродинамики. Выпускная
квалификационная работа бакалавра. МГТУ им.~Н.\,Э.~Баумана, 2025.
\bibitem{toth2000} T\'oth~G. The $\nabla\!\cdot\!\B=0$ constraint in
shock-capturing magnetohydrodynamics codes // J. Comput. Phys. 2000. Vol.~161.
P.~605--652.
\bibitem{balsara1999} Balsara~D.\,S., Spicer~D.\,S. A staggered mesh algorithm
using high order Godunov fluxes to ensure solenoidal magnetic fields in MHD
simulations // J. Comput. Phys. 1999. Vol.~149. P.~270--292.
\bibitem{gardiner2005} Gardiner~T.\,A., Stone~J.\,M. An unsplit Godunov method
for ideal MHD via constrained transport // J. Comput. Phys. 2005. Vol.~205.
P.~509--539.
\bibitem{shu1988} Shu~C.-W., Osher~S. Efficient implementation of essentially
non-oscillatory shock-capturing schemes // J. Comput. Phys. 1988. Vol.~77.
P.~439--471.
\bibitem{gottlieb2001} Gottlieb~S., Shu~C.-W., Tadmor~E. Strong
stability-preserving high-order time discretization methods // SIAM Review.
2001. Vol.~43. P.~89--112.
\bibitem{amrex} Zhang~W. et al. AMReX: a framework for block-structured
adaptive mesh refinement // J. Open Source Softw. 2019. Vol.~4, № 37. 1370.
\bibitem{toro2009} Toro~E.\,F. Riemann Solvers and Numerical Methods for Fluid
Dynamics: A Practical Introduction. 3rd ed. Berlin: Springer, 2009. 721~p.
\bibitem{takahashi2013} Takahashi~K., Yamada~S. Regular and non-regular
solutions of the Riemann problem in ideal magnetohydrodynamics //
J. Plasma Phys. 2013. Vol.~79. P.~335--356 (arXiv:1210.5584).
\bibitem{zhilkin2010} Жилкин~А.\,Г., Бисикало~Д.\,В. Формирование аккреционных
дисков в тесных двойных системах с магнитными полями // Astronomy Reports. 2010.
Vol.~54, № 12. P.~1063. arXiv:1010.4949.
\bibitem{zhilkin2012} Жилкин~А.\,Г., Бисикало~Д.\,В., Боярчук~А.\,А. Структура
течения в тесных двойных звёздах с учётом магнитного поля // УФН. 2012. Т.~182,
№ 2. С.~121--145. DOI 10.3367/UFNr.0182.201202a.0121.
\bibitem{bisikalo2013} Бисикало~Д.\,В., Жилкин~А.\,Г., Боярчук~А.\,А.
Газодинамика тесных двойных звёзд. М.: Физматлит, 2013. 632~с.
\bibitem{dudorov1999} Дудоров~А.\,Е., Жилкин~А.\,Г., Кузнецов~О.\,А.
Квазимонотонная разностная схема повышенного порядка точности для уравнений
магнитной гидродинамики // Матем. моделирование. 1999. Т.~11, № 1.
\bibitem{kulikovskii2005} Куликовский~А.\,Г., Любимов~Г.\,А. Магнитная
гидродинамика. М.: Логос, 2005. 328~с.
\end{thebibliography}
\end{document}
